%%% Title presen22101
\documentclass[landscape,10pt]{ujarticle}
\special{papersize=\the\paperwidth,\the\paperheight}
\usepackage{ketpic,ketlayer}
\usepackage{ketslide}
\usepackage{amsmath,amssymb}
\usepackage{bm,enumerate}
\usepackage[dvipdfmx]{graphicx}
\usepackage{color}
\definecolor{slidecolora}{cmyk}{0.98,0.13,0,0.43}
\definecolor{slidecolorb}{cmyk}{0.2,0,0,0}
\definecolor{slidecolorc}{cmyk}{0.2,0,0,0}
\definecolor{slidecolord}{cmyk}{0.2,0,0,0}
\definecolor{slidecolore}{cmyk}{0,0,0,0.5}
\definecolor{slidecolorf}{cmyk}{0,0,0,0.5}
\definecolor{slidecolori}{cmyk}{0.98,0.13,0,0.43}
\def\setthin#1{\def\thin{#1}}
\setthin{0}
\newcommand{\slidepage}[1][s]{%
\setcounter{ketpicctra}{18}%
\if#1m \setcounter{ketpicctra}{1}\fi
\hypersetup{linkcolor=black}%

\begin{layer}{118}{0}
\putnotee{122}{-\theketpicctra.05}{\small\thepage/\pageref{pageend}}
\end{layer}\hypersetup{linkcolor=blue}

}
\usepackage{pict2e}
\usepackage[dvipdfmx,colorlinks=true,linkcolor=blue,filecolor=blue]{hyperref}
\newcommand{\hiduke}{0411}
\newcommand{\hako}[2][1]{\fbox{\raisebox{#1mm}{\mbox{}}\raisebox{-#1mm}{\mbox{}}\,\phantom{#2}\,}}
\newcommand{\hakoa}[2][1]{\fbox{\raisebox{#1mm}{\mbox{}}\raisebox{-#1mm}{\mbox{}}\,#2\,}}
\newcommand{\hakom}[2][1]{\hako[#1]{$#2$}}
\newcommand{\hakoma}[2][1]{\hakoa[#1]{$#2$}}
\def\rad{\;\mathrm{rad}}
\def\deg#1{#1^{\circ}}
\newcommand{\sbunsuu}[2]{\scalebox{0.6}{$\bunsuu{#1}{#2}$}}
\def\pow{$\hspace{-1.5mm}＾\hspace{-1mm}$}
\def\dlim{\displaystyle\lim}
\newcommand{\brd}[2][1]{\scalebox{#1}{\color{red}\fbox{\color{black}$#2$}}}

\setmargin{25}{145}{15}{100}

\ketslideinit

\pagestyle{empty}

\begin{document}

\begin{layer}{120}{0}
\putnotese{0}{0}{{\Large\bf
\color[cmyk]{1,1,0,0}

\begin{layer}{120}{0}
{\Huge \putnotes{60}{20}{この授業について}}
\putnotes{60}{50}{高遠節夫(たかとおせつお)}
\putnotes{60}{70}{2022.04.11}
\end{layer}

}
}
\end{layer}

\def\mainslidetitley{22}
\def\ketcletter{slidecolora}
\def\ketcbox{slidecolorb}
\def\ketdbox{slidecolorc}
\def\ketcframe{slidecolord}
\def\ketcshadow{slidecolore}
\def\ketdshadow{slidecolorf}
\def\slidetitlex{6}
\def\slidetitlesize{1.3}
\def\mketcletter{slidecolori}
\def\mketcbox{yellow}
\def\mketdbox{yellow}
\def\mketcframe{yellow}
\def\mslidetitlex{62}
\def\mslidetitlesize{2}

\color{black}
\Large\bf\boldmath
\addtocounter{page}{-1}

\def\MARU{}
\renewcommand{\MARU}[1]{{\ooalign{\hfil$#1$\/\hfil\crcr\raise.167ex\hbox{\mathhexbox20D}}}}
\renewcommand{\slidepage}[1][s]{%
\setcounter{ketpicctra}{18}%
\if#1m \setcounter{ketpicctra}{1}\fi
\hypersetup{linkcolor=black}%
\begin{layer}{118}{0}
\putnotee{115}{-\theketpicctra.05}{\small\hiduke-\thepage/\pageref{pageend}}
\end{layer}\hypersetup{linkcolor=blue}
}
\newcounter{ban}
\setcounter{ban}{1}
\newcommand{\monban}[1][\hiduke]{%
#1-\theban\ %
\addtocounter{ban}{1}%
}
\newcommand{\monbannoadd}[1][\hiduke]{%
#1-\theban\ %
}
\newcommand{\addban}{%
\addtocounter{ban}{1}%%210614
}
\newcounter{edawidth}
\newcounter{edactr}
\newcommand{\seteda}[1]{
\setcounter{edawidth}{#1}
\setcounter{edactr}{1}
}
\newcommand{\eda}[2][\theedawidth ]{%
\noindent\Ltab{#1 mm}{[\theedactr]\ #2}%
\addtocounter{edactr}{1}%
}
%%%%%%%%%%%%

%%%%%%%%%%%%%%%%%%%%

\mainslide{Google Classroom(GC)}


\slidepage[m]
%%%%%%%%%%%%

%%%%%%%%%%%%%%%%%%%%

\newslide{準備}

\vspace*{18mm}

\slidepage
\begin{enumerate}[(1)]
\item
Gmailのアカウントがない場合は作成する
\item
スマホ：GoogleClassroomのアプリを入手（無料）\\
PC　：「GoogleClassroom ログイン」で検索
\item
\verb|Google Classroom|(以下GC)にログインする\\
\hspace*{1zw}注）アカウント名とパスワードが必要
\item
右上か右下の「＋」を押して「クラスに参加」を選ぶ\\
\hspace*{1zw}クラスコード {\color{red}6mxdy4j} を入力
\end{enumerate}
%%%%%%%%%%%%

%%%%%%%%%%%%%%%%%%%%

\newslide{授業(polytech22)のページ}

\vspace*{18mm}

\slidepage

\begin{layer}{120}{0}
\putnotes{65}{2}{\includegraphics[bb=0.00 0.00 2210.00 1540.00,width=90mm]{fig/classroomtop.pdf}}
\end{layer}

%%%%%%%%%%%%

%%%%%%%%%%%%%%%%%%%%


\newslide{授業の方法}

\vspace*{18mm}

%%repeat=4
\slidepage
\begin{itemize}
\item
主にスライドを用いる
\item
ノートに要点をまとめ，問題を解く
\item
試験ではノートの持ち込みOK
\item
教材や課題の配付は，GCから%%\\
%%\hspace*{1zw}HTML教材のときは，スマホを横置きで使う
\item
{\color{red}授業中は，指示した時以外スマホを使わないこと}
\end{itemize}
%%%%%%%%%%%%

%%%%%%%%%%%%%%%%%%%%

\mainslide{数式のやりとり}


\slidepage[m]
%%%%%%%%%%%%

%%%%%%%%%%%%%%%%%%%%

\newslide{KeTMathの利用}

\vspace*{18mm}

\slidepage
\begin{itemize}
\item
普通の数式(2次元記法)は見やすい．\\
　$\dfrac{4}{9}$,\ $\sqrt{7}$,\ $5^3$
\item
しかし，オンラインでのやりとりには向かない\\
　＝＞1次元記法の方がよい
\item
KeTMath\\
　・1次元数式を入力すると即時に2次元数式を表示\\
　・問題の1次元数式を入力窓1にコピーして確認\\
　・解答を入力窓2に入力してよければGCにコピー
\end{itemize}

\newslide{1次元数式の記法}

\vspace*{18mm}

%%repeat=4
\slidepage
\begin{itemize}
\item
\Ltab{8zw}{分数 (fraction)}$\dfrac{a}{b}\ \Longrightarrow$\ \verb|fr(a,b)|
\item
\Ltab{8zw}{割り算}$a\div b\ \Longrightarrow$\ \verb|a {\div} b|\\
\hfill（あまり使わない）
\item
\Ltab{6zw}{掛け算}$ab\ \Longrightarrow$\ \verb|ab, a*b, a(cdot)b, a(cross)b|
\item
\Ltab{10zw}{べき乗}$a^b\ \Longrightarrow$\ \verb|a^(b)|
\item
\Ltab{12zw}{平方根 (square root)}$\sqrt{a}\ \Longrightarrow$\ \verb|sq(a)|
\item
\Ltab{10zw}{円周率}$\pi\ \Longrightarrow$\ \verb|pi|
\end{itemize}
%%%%%%%%%%%%

%%%%%%%%%%%%%%%%%%%%

\newslide{GCとKeTMath}

\vspace*{18mm}

%%repeat=2
\slidepage
\begin{enumerate}[(1)]
\item
GCの質問にリンクがあるときはリンクをクリック\vspace{-2mm}
\item
課題を埋め込んだKeTMathが立ち上がる\vspace{-2mm}
\item
自分の番号を入れてOKを押す\vspace{-2mm}
\item
入力2に解答を入れる\vspace{-2mm}
\item
Recを押すと全ての解答が入力3に入る\vspace{-2mm}
\item
入力3の「すべてを選択」コピーする\vspace{-2mm}
\item
GCの回答欄にペーストして送信を押す
\end{enumerate}
%%%%%%%%%%%%

%%%%%%%%%%%%%%%%%%%%

\newslide{課題(KeTMathに慣れよう)}

\vspace*{18mm}

%%repeat=2
\slidepage
\begin{itemize}
\item
課題 \monban 次の1次元数式の2次元数式を書け．\seteda{55}\\
\eda{fr(1+4,3)}\eda{a+b/c+d}\\
\eda{3sq(6)}\eda{pir\^{}\!(2)}
\item
課題 \monban 次の数式を1次元数式で書け．\vspace{2mm}\seteda{55}\\
\eda{$-\dfrac{3}{5}$}\eda{$\dfrac{xy}{x+y}$}\vspace{2mm}\\
\eda{$\sqrt{3}-\sqrt{2}$}\eda{$\dfrac{\pi}{2}$}
\end{itemize}
%%%%%%%%%%%%

%%%%%%%%%%%%%%%%%%%%

\mainslide{数と式の計算}


\slidepage[m]
%%%%%%%%%%%%

%%%%%%%%%%%%%%%%%%%%

\newslide{正負の数の足し算と引き算}

\vspace*{18mm}

\slidepage
\begin{itemize}
\item
正の数$-$正の数\\
\hspace*{2zw}$14-6=8$\\
\hspace*{2zw}$6-8=-2$
\item
正の数$+$負の数\\
\hspace*{2zw}$12+(-3)=12-3=9$\\
\hspace*{2zw}$(-5)+3=3+(-5)=3-5=-2$
\item
負の数$+$負の数\\
\hspace*{2zw}$(-2)+(-3)=-(2+3)=-5$
\end{itemize}

\newslide{正負の数の掛け算（割り算）}

\vspace*{18mm}

\slidepage
\begin{itemize}
\item
正の数$\times$負の数$=$負の数\\
\hspace*{2zw}$6\times(-3)=-18$
\item
負の数$\times$負の数$=$正の数\\
\hspace*{2zw}$(-4)\times(-3)=12$\\
\hspace*{1zw}注）$(-4)(-3)$とか$(-4)\cdot(-3)$と書くこともある
\end{itemize}

\newslide{計算問題}

\vspace*{18mm}

%%repeat=2
\slidepage
\seteda{55}
\begin{itemize}
\item
課題\monban\\
\eda{$-6+5$}\eda{$8-(-2)$}\\
\eda{$(-7)(+8)$}\eda{$32\div(-4)\times 8$}
\seteda{55}
\item
課題\monban\\
\eda{$6-8\div (-4)$}\eda{$18\div(-6)-7\times(-2)$}\\
\eda{$54\div(3^2-3)$}\eda{$3\times 23 + 3\times 77$}
\end{itemize}
%%%%%%%%%%%%

%%%%%%%%%%%%%%%%%%%%

\newslide{分数の計算}

\vspace*{18mm}

\slidepage
\begin{itemize}
\item
約分　分母と分子を同じ数で割る\vspace{1mm}\\
\hspace*{2zw}$\dfrac{4}{6}=\dfrac{2}{3}$（分母と分子を$2$で割る）
\item
通分　2つの分数の分母を同じにする
\item
足し算（引き算）　通分してから分子どうしを計算\vspace{1mm}\\
\hspace*{2zw}$\dfrac{3}{4}+\dfrac{1}{6}=\dfrac{9}{12}+\dfrac{2}{12}=\dfrac{9+2}{12}=\dfrac{11}{12}$
\item
掛け算　分母どうし，分子どうしを掛ける\vspace{1mm}\\
\hspace*{2zw}$\dfrac{2}{5}\times\dfrac{5}{6}=\dfrac{10}{30}=\dfrac{1}{3}$
\end{itemize}

\newslide{分数の計算問題}

\vspace*{18mm}

%%repeat=2,para
\slidepage
\seteda{55}
\begin{itemize}
\item
課題\monban\vspace{1mm}\\
\eda{$\dfrac{1}{2}-\dfrac{1}{4}$}\eda{$\dfrac{2}{3}-\dfrac{5}{6}$}\vspace{1mm}\\
\eda{$\dfrac{3}{10}-\dfrac{3}{5}$}\eda{$\dfrac{1}{2}+\dfrac{1}{3}$}
\item
課題\monban\vspace{1mm}\seteda{55}\\
\eda{$\dfrac{4}{5}\times\dfrac{2}{3}$}\eda{$\dfrac{2}{5}\times\dfrac{3}{7}$}\vspace{1mm}\\
\eda{$\dfrac{4}{3}\times\dfrac{9}{8}$}\eda{$\dfrac{4}{9}\times\dfrac{6}{5}$}
\end{itemize}
%%%%%%%%%%%%

%%%%%%%%%%%%%%%%%%%%

\newslide{分数の割り算}

\vspace*{18mm}

%%repeat=1
\slidepage
\seteda{55}
\begin{itemize}
\item
割る方の分母と分子をひっくり返して掛ける\vspace{1mm}\\
\hspace*{2zw}$\dfrac{9}{26}\div\dfrac{3}{4}=\dfrac{9}{26}\times\dfrac{4}{3}=
\dfrac{9\times 4}{26\times 3}=\dfrac{3 \times 2}{13}=\dfrac{6}{13}$
\item
課題\monban\vspace{1mm}\\
\eda{$\dfrac{4}{7}\div\dfrac{2}{3}$}\eda{$\dfrac{2}{5}\div\dfrac{4}{7}$}\vspace{1mm}\\
\eda{$\dfrac{7}{12}\div\dfrac{3}{8}$}\eda{$\dfrac{5}{4}\div\dfrac{15}{7}$}
\end{itemize}
%%%%%%%%%%%%

%%%%%%%%%%%%%%%%%%%%

\newslide{文字式}

\vspace*{18mm}

\slidepage

\begin{layer}{120}{0}
\putnotese{70}{14}{%%% /Users/takatoosetsuo/Dropbox/2022polytech/101/fig/presen22101rope.tex 
%%% Generator=presen22101.cdy 
{\unitlength=1cm%
\begin{picture}%
(4,0.8)(0,0)%
\linethickness{0.008in}%%
\Large\bf\boldmath%
\small%
\linethickness{0.024in}%%
\polyline(0.00000,0.25000)(4.00000,0.25000)%
%
\linethickness{0.008in}%%
\polyline(4.00000,0.25000)(3.96322,0.26196)(3.92636,0.27369)(3.88943,0.28521)(3.85244,0.29650)%
(3.81538,0.30756)(3.77825,0.31840)(3.74105,0.32902)(3.70380,0.33942)(3.66648,0.34958)%
(3.62910,0.35953)(3.59166,0.36924)(3.55417,0.37874)(3.51661,0.38800)(3.47900,0.39704)%
(3.44134,0.40585)(3.40363,0.41443)(3.36586,0.42279)(3.32805,0.43092)(3.29018,0.43882)%
(3.25227,0.44649)(3.21431,0.45393)(3.17631,0.46114)(3.13827,0.46813)(3.10018,0.47488)%
(3.06206,0.48140)(3.02390,0.48770)(2.98570,0.49376)(2.94746,0.49959)(2.90919,0.50519)%
(2.87088,0.51056)(2.83255,0.51570)(2.79418,0.52061)(2.75578,0.52529)(2.71736,0.52973)%
(2.67891,0.53395)(2.64044,0.53793)(2.60194,0.54167)(2.56342,0.54519)(2.52488,0.54847)%
(2.48632,0.55152)(2.44775,0.55434)(2.40916,0.55693)(2.37055,0.55928)(2.33193,0.56140)%
(2.29329,0.56329)(2.25465,0.56494)(2.21600,0.56636)(2.17734,0.56755)(2.13867,0.56850)%
(2.10000,0.56922)%
%
\polyline(1.90000,0.56922)(1.86133,0.56850)(1.82266,0.56755)(1.78400,0.56636)(1.74535,0.56494)%
(1.70671,0.56329)(1.66807,0.56140)(1.62945,0.55928)(1.59084,0.55693)(1.55225,0.55434)%
(1.51368,0.55152)(1.47512,0.54847)(1.43658,0.54519)(1.39806,0.54167)(1.35956,0.53793)%
(1.32109,0.53395)(1.28264,0.52973)(1.24422,0.52529)(1.20582,0.52061)(1.16745,0.51570)%
(1.12912,0.51056)(1.09081,0.50519)(1.05254,0.49959)(1.01430,0.49376)(0.97610,0.48770)%
(0.93794,0.48140)(0.89982,0.47488)(0.86173,0.46813)(0.82369,0.46114)(0.78569,0.45393)%
(0.74773,0.44649)(0.70982,0.43882)(0.67195,0.43092)(0.63414,0.42279)(0.59637,0.41443)%
(0.55866,0.40585)(0.52100,0.39704)(0.48339,0.38800)(0.44583,0.37874)(0.40834,0.36924)%
(0.37090,0.35953)(0.33352,0.34958)(0.29620,0.33942)(0.25895,0.32902)(0.22175,0.31840)%
(0.18462,0.30756)(0.14756,0.29650)(0.11057,0.28521)(0.07364,0.27369)(0.03678,0.26196)%
(-0.00000,0.25000)%
%
\settowidth{\Width}{$a$}\setlength{\Width}{-0.5\Width}%
\settoheight{\Height}{$a$}\settodepth{\Depth}{$a$}\setlength{\Height}{-0.5\Height}\setlength{\Depth}{0.5\Depth}\addtolength{\Height}{\Depth}%
\put(2.0000000,0.5700000){\hspace*{\Width}\raisebox{\Height}{$a$}}%
%
\polyline(2.50000,0.25000)(2.51238,0.24489)(2.52480,0.23986)(2.53726,0.23493)(2.54974,0.23007)%
(2.56226,0.22530)(2.57481,0.22062)(2.58740,0.21603)(2.60001,0.21152)(2.61266,0.20710)%
(2.62534,0.20276)(2.63804,0.19851)(2.65078,0.19435)(2.66354,0.19028)(2.67633,0.18629)%
(2.68915,0.18240)(2.70199,0.17859)(2.71486,0.17487)(2.72775,0.17123)(2.74067,0.16769)%
(2.75362,0.16424)(2.76658,0.16087)(2.77957,0.15759)(2.79259,0.15441)(2.80562,0.15131)%
(2.81868,0.14830)(2.83175,0.14538)(2.84485,0.14255)(2.85796,0.13982)(2.87109,0.13717)%
(2.88424,0.13461)(2.89741,0.13214)(2.91059,0.12976)(2.92379,0.12748)(2.93701,0.12528)%
(2.95024,0.12318)(2.96349,0.12116)(2.97674,0.11924)(2.99001,0.11741)(3.00330,0.11567)%
(3.01659,0.11402)(3.02990,0.11246)(3.04321,0.11100)(3.05654,0.10962)(3.06988,0.10834)%
(3.08322,0.10715)(3.09657,0.10605)(3.10993,0.10504)(3.12330,0.10412)(3.13667,0.10330)%
(3.15004,0.10256)%
%
\polyline(3.34996,0.10256)(3.36333,0.10330)(3.37670,0.10412)(3.39007,0.10504)(3.40343,0.10605)%
(3.41678,0.10715)(3.43012,0.10834)(3.44346,0.10962)(3.45679,0.11100)(3.47010,0.11246)%
(3.48341,0.11402)(3.49670,0.11567)(3.50999,0.11741)(3.52326,0.11924)(3.53651,0.12116)%
(3.54976,0.12318)(3.56299,0.12528)(3.57621,0.12748)(3.58941,0.12976)(3.60259,0.13214)%
(3.61576,0.13461)(3.62891,0.13717)(3.64204,0.13982)(3.65515,0.14255)(3.66825,0.14538)%
(3.68132,0.14830)(3.69438,0.15131)(3.70741,0.15441)(3.72043,0.15759)(3.73342,0.16087)%
(3.74638,0.16424)(3.75933,0.16769)(3.77225,0.17123)(3.78514,0.17487)(3.79801,0.17859)%
(3.81085,0.18240)(3.82367,0.18629)(3.83646,0.19028)(3.84922,0.19435)(3.86196,0.19851)%
(3.87466,0.20276)(3.88734,0.20710)(3.89999,0.21152)(3.91260,0.21603)(3.92519,0.22062)%
(3.93774,0.22530)(3.95026,0.23007)(3.96274,0.23493)(3.97520,0.23986)(3.98762,0.24489)%
(4.00000,0.25000)%
%
\settowidth{\Width}{$x$}\setlength{\Width}{-0.5\Width}%
\settoheight{\Height}{$x$}\settodepth{\Depth}{$x$}\setlength{\Height}{-0.5\Height}\setlength{\Depth}{0.5\Depth}\addtolength{\Height}{\Depth}%
\put(3.2500000,0.1000000){\hspace*{\Width}\raisebox{\Height}{$x$}}%
%
\linethickness{0.004in}%%
\polyline(-0.00000,0.25000)(0.02148,0.24392)(0.04299,0.23795)(0.06453,0.23208)(0.08610,0.22632)%
(0.10770,0.22066)(0.12933,0.21511)(0.15098,0.20967)(0.17265,0.20433)(0.19436,0.19910)%
(0.21609,0.19398)(0.23784,0.18896)(0.25962,0.18405)(0.28142,0.17925)(0.30325,0.17455)%
(0.32510,0.16996)(0.34697,0.16548)(0.36886,0.16110)(0.39077,0.15683)(0.41271,0.15267)%
(0.43466,0.14862)(0.45664,0.14467)(0.47863,0.14084)(0.50064,0.13711)(0.52267,0.13349)%
(0.54472,0.12997)(0.56678,0.12657)(0.58886,0.12327)(0.61096,0.12008)(0.63307,0.11700)%
(0.65520,0.11403)(0.67734,0.11117)(0.69949,0.10841)(0.72166,0.10577)(0.74384,0.10323)%
(0.76604,0.10080)(0.78824,0.09848)(0.81046,0.09627)(0.83268,0.09417)(0.85492,0.09218)%
(0.87716,0.09030)(0.89942,0.08852)(0.92168,0.08686)(0.94395,0.08530)(0.96623,0.08385)%
(0.98852,0.08252)(1.01081,0.08129)(1.03311,0.08017)(1.05541,0.07916)(1.07772,0.07826)%
(1.10003,0.07747)%
%
\polyline(1.39997,0.07747)(1.42228,0.07826)(1.44459,0.07916)(1.46689,0.08017)(1.48919,0.08129)%
(1.51148,0.08252)(1.53377,0.08385)(1.55605,0.08530)(1.57832,0.08686)(1.60058,0.08852)%
(1.62284,0.09030)(1.64508,0.09218)(1.66732,0.09417)(1.68954,0.09627)(1.71176,0.09848)%
(1.73396,0.10080)(1.75616,0.10323)(1.77834,0.10577)(1.80051,0.10841)(1.82266,0.11117)%
(1.84480,0.11403)(1.86693,0.11700)(1.88904,0.12008)(1.91114,0.12327)(1.93322,0.12657)%
(1.95528,0.12997)(1.97733,0.13349)(1.99936,0.13711)(2.02137,0.14084)(2.04336,0.14467)%
(2.06534,0.14862)(2.08729,0.15267)(2.10923,0.15683)(2.13114,0.16110)(2.15303,0.16548)%
(2.17490,0.16996)(2.19675,0.17455)(2.21858,0.17925)(2.24038,0.18405)(2.26216,0.18896)%
(2.28391,0.19398)(2.30564,0.19910)(2.32735,0.20433)(2.34902,0.20967)(2.37067,0.21511)%
(2.39230,0.22066)(2.41390,0.22632)(2.43547,0.23208)(2.45701,0.23795)(2.47852,0.24392)%
(2.50000,0.25000)%
%
\linethickness{0.008in}%%
\settowidth{\Width}{$b$}\setlength{\Width}{-0.5\Width}%
\settoheight{\Height}{$b$}\settodepth{\Depth}{$b$}\setlength{\Height}{-0.5\Height}\setlength{\Depth}{0.5\Depth}\addtolength{\Height}{\Depth}%
\put(1.2500000,0.0800000){\hspace*{\Width}\raisebox{\Height}{$b$}}%
%
\end{picture}}%}
\end{layer}

\seteda{110}
\begin{itemize}
\item
長さ$a$のひもから長さ$x$のひもを切り取ったときの残りの長さを$b$とする．\\
\hspace*{4zw}$b=a-x$
\item
課題\monbannoadd 次を文字式で表せ\hfill(TextP4)\\
\eda{1辺の長さが$a$\;cmである正方形の面積$S=$}\\
\eda{円周率$\pi$，半径が$r$である円の円周$L=$}
\end{itemize}
\addban

\newslide{文字式の計算}

\vspace*{18mm}

\slidepage
\begin{itemize}
\item
掛け算記号は省略　$x\cdot y,\ x\times y\ \Longrightarrow\ xy$
\item
べき乗　$xx,\ xxx\ \Longrightarrow\ x^2,\ x^3$
\item
数は文字の前におく　$x\cdot 3\cdot y\cdot 4=12xy$
\item
計算は，数の場合と同様\\
\hspace*{2zw}$3a\times (-7a^2)=3\cdot(-7) aa^2=-21 a^3$
\item
課題\monbannoadd 次の計算をせよ．\hfill(TextP5)\seteda{55}\\
\eda{$-\dfrac{9}{2}a\times\left(-\dfrac{5}{6}b\right)$}%
\eda{$\dfrac{2}{3}a\times(-3a)^2$}
\end{itemize}
\addban

\mainslide{関数}


\slidepage[m]
%%%%%%%%%%%%

%%%%%%%%%%%%%%%%%%%%

\newslide{関数}

\vspace*{18mm}

\slidepage
\begin{itemize}
\item
変数$x$の値を与えると変数$y$の値が求まる\\
\hspace*{1zw}例）$y=2x+1,\ y=x^2+2x+1$
\item
これを変数$x$の関数という
\item
変数$x$の関数であることを$f(x)$などで表す\\
　例1）$f(x)=2x+1$（1次関数）\\
　例2）$g(x)=x^2+2x+1$（2次関数）
\end{itemize}

\newslide{関数記号}

\vspace*{18mm}

\slidepage
\seteda{55}
\begin{itemize}
\item
関数$f(x)$の$x$に定数$a$を代入した値を$f(a)$で表す
\item
例）$f(x)=x^2+x-1$のとき\\
　　　$f(2)=2^2+2-1=5$
\item
課題\monban $f(x)=3x+1$のとき，次を求めよ．\\
\eda{$f(0)$}\eda{$f(2)$}\\
\eda{$f(-3)$}\eda{$f(a-1)$（$a$は定数）}\\
\hfill(TextP80問1,2)\\
\end{itemize}

\newslide{関数のグラフ}

\vspace*{18mm}

%%repeat=8,para
\slidepage
\vspace{5mm}

関数$y=f(x)$
\begin{itemize}
\item
$x$を変えるとき，点$\bigl(x,\ f(x)\bigr)$も変わる．
\item
[]例) 1次関数$y=2x+1$\vspace{1mm}\\
%%% /polytech22.git/102(0418)/presen/fig/table1a.tex 
%%% Generator=presen22102.cdy 
{\unitlength=1cm%
\begin{picture}%
(9.6,1.2)(0,0)%
\linethickness{0.008in}%%
\Large\bf\boldmath%
\small%
\polyline(0.00000,1.20000)(0.00000,0.00000)%
%
\polyline(0.80000,1.20000)(0.80000,0.00000)%
%
\polyline(1.60000,1.20000)(1.60000,0.00000)%
%
\polyline(2.40000,1.20000)(2.40000,0.00000)%
%
\polyline(3.20000,1.20000)(3.20000,0.00000)%
%
\polyline(4.00000,1.20000)(4.00000,0.00000)%
%
\polyline(4.80000,1.20000)(4.80000,0.00000)%
%
\polyline(5.60000,1.20000)(5.60000,0.00000)%
%
\polyline(6.40000,1.20000)(6.40000,0.00000)%
%
\polyline(7.20000,1.20000)(7.20000,0.00000)%
%
\polyline(8.00000,1.20000)(8.00000,0.00000)%
%
\polyline(8.80000,1.20000)(8.80000,0.00000)%
%
\polyline(9.60000,1.20000)(9.60000,0.00000)%
%
\polyline(0.00000,1.20000)(9.60000,1.20000)%
%
\polyline(0.00000,0.60000)(9.60000,0.60000)%
%
\polyline(0.00000,0.00000)(9.60000,0.00000)%
%
\settowidth{\Width}{$x$}\setlength{\Width}{-0.5\Width}%
\settoheight{\Height}{$x$}\settodepth{\Depth}{$x$}\setlength{\Height}{-0.5\Height}\setlength{\Depth}{0.5\Depth}\addtolength{\Height}{\Depth}%
\put(0.4000000,0.9000000){\hspace*{\Width}\raisebox{\Height}{$x$}}%
%
\settowidth{\Width}{$-5$}\setlength{\Width}{-0.5\Width}%
\settoheight{\Height}{$-5$}\settodepth{\Depth}{$-5$}\setlength{\Height}{-0.5\Height}\setlength{\Depth}{0.5\Depth}\addtolength{\Height}{\Depth}%
\put(1.2000000,0.9000000){\hspace*{\Width}\raisebox{\Height}{$-5$}}%
%
\settowidth{\Width}{$-4$}\setlength{\Width}{-0.5\Width}%
\settoheight{\Height}{$-4$}\settodepth{\Depth}{$-4$}\setlength{\Height}{-0.5\Height}\setlength{\Depth}{0.5\Depth}\addtolength{\Height}{\Depth}%
\put(2.0000000,0.9000000){\hspace*{\Width}\raisebox{\Height}{$-4$}}%
%
\settowidth{\Width}{$-3$}\setlength{\Width}{-0.5\Width}%
\settoheight{\Height}{$-3$}\settodepth{\Depth}{$-3$}\setlength{\Height}{-0.5\Height}\setlength{\Depth}{0.5\Depth}\addtolength{\Height}{\Depth}%
\put(2.8000000,0.9000000){\hspace*{\Width}\raisebox{\Height}{$-3$}}%
%
\settowidth{\Width}{$-2$}\setlength{\Width}{-0.5\Width}%
\settoheight{\Height}{$-2$}\settodepth{\Depth}{$-2$}\setlength{\Height}{-0.5\Height}\setlength{\Depth}{0.5\Depth}\addtolength{\Height}{\Depth}%
\put(3.6000000,0.9000000){\hspace*{\Width}\raisebox{\Height}{$-2$}}%
%
\settowidth{\Width}{$-1$}\setlength{\Width}{-0.5\Width}%
\settoheight{\Height}{$-1$}\settodepth{\Depth}{$-1$}\setlength{\Height}{-0.5\Height}\setlength{\Depth}{0.5\Depth}\addtolength{\Height}{\Depth}%
\put(4.4000000,0.9000000){\hspace*{\Width}\raisebox{\Height}{$-1$}}%
%
\settowidth{\Width}{$0$}\setlength{\Width}{-0.5\Width}%
\settoheight{\Height}{$0$}\settodepth{\Depth}{$0$}\setlength{\Height}{-0.5\Height}\setlength{\Depth}{0.5\Depth}\addtolength{\Height}{\Depth}%
\put(5.2000000,0.9000000){\hspace*{\Width}\raisebox{\Height}{$0$}}%
%
\settowidth{\Width}{$1$}\setlength{\Width}{-0.5\Width}%
\settoheight{\Height}{$1$}\settodepth{\Depth}{$1$}\setlength{\Height}{-0.5\Height}\setlength{\Depth}{0.5\Depth}\addtolength{\Height}{\Depth}%
\put(6.0000000,0.9000000){\hspace*{\Width}\raisebox{\Height}{$1$}}%
%
\settowidth{\Width}{$2$}\setlength{\Width}{-0.5\Width}%
\settoheight{\Height}{$2$}\settodepth{\Depth}{$2$}\setlength{\Height}{-0.5\Height}\setlength{\Depth}{0.5\Depth}\addtolength{\Height}{\Depth}%
\put(6.8000000,0.9000000){\hspace*{\Width}\raisebox{\Height}{$2$}}%
%
\settowidth{\Width}{$3$}\setlength{\Width}{-0.5\Width}%
\settoheight{\Height}{$3$}\settodepth{\Depth}{$3$}\setlength{\Height}{-0.5\Height}\setlength{\Depth}{0.5\Depth}\addtolength{\Height}{\Depth}%
\put(7.6000000,0.9000000){\hspace*{\Width}\raisebox{\Height}{$3$}}%
%
\settowidth{\Width}{$4$}\setlength{\Width}{-0.5\Width}%
\settoheight{\Height}{$4$}\settodepth{\Depth}{$4$}\setlength{\Height}{-0.5\Height}\setlength{\Depth}{0.5\Depth}\addtolength{\Height}{\Depth}%
\put(8.4000000,0.9000000){\hspace*{\Width}\raisebox{\Height}{$4$}}%
%
\settowidth{\Width}{$5$}\setlength{\Width}{-0.5\Width}%
\settoheight{\Height}{$5$}\settodepth{\Depth}{$5$}\setlength{\Height}{-0.5\Height}\setlength{\Depth}{0.5\Depth}\addtolength{\Height}{\Depth}%
\put(9.2000000,0.9000000){\hspace*{\Width}\raisebox{\Height}{$5$}}%
%
\settowidth{\Width}{$y$}\setlength{\Width}{-0.5\Width}%
\settoheight{\Height}{$y$}\settodepth{\Depth}{$y$}\setlength{\Height}{-0.5\Height}\setlength{\Depth}{0.5\Depth}\addtolength{\Height}{\Depth}%
\put(0.4000000,0.3000000){\hspace*{\Width}\raisebox{\Height}{$y$}}%
%
\end{picture}}%
\item
この点の集まりを，その関数の{\color{red}グラフ}という．
\end{itemize}
\if 1=0 % to end
%%%%%%%%%%%%

%%%%%%%%%%%%%%%%%%%%

\newslide{1次関数のグラフ}

\vspace*{18mm}

\slidepage
例）$y=2x+1$

\begin{layer}{120}{0}
\putnotese{70}{-3}{\scalebox{0.5}{%%% /polytech.git/n101/fig/table1b.tex 
%%% Generator=presen0601.cdy 
{\unitlength=1cm%
\begin{picture}%
(9.6,1.2)(0,0)%
\special{pn 8}%
%
\Large\bf\boldmath%
\small%
{%
\color{red}%
\special{pa     0  -472}\special{pa     0    -0}%
\special{fp}%
}%
{%
\color{red}%
\special{pa   315  -472}\special{pa   315    -0}%
\special{fp}%
}%
{%
\color{red}%
\special{pa   630  -472}\special{pa   630    -0}%
\special{fp}%
}%
{%
\color{red}%
\special{pa   945  -472}\special{pa   945    -0}%
\special{fp}%
}%
{%
\color{red}%
\special{pa  1260  -472}\special{pa  1260    -0}%
\special{fp}%
}%
{%
\color{red}%
\special{pa  1575  -472}\special{pa  1575    -0}%
\special{fp}%
}%
{%
\color{red}%
\special{pa  1890  -472}\special{pa  1890    -0}%
\special{fp}%
}%
{%
\color{red}%
\special{pa  2205  -472}\special{pa  2205    -0}%
\special{fp}%
}%
{%
\color{red}%
\special{pa  2520  -472}\special{pa  2520    -0}%
\special{fp}%
}%
{%
\color{red}%
\special{pa  2835  -472}\special{pa  2835    -0}%
\special{fp}%
}%
{%
\color{red}%
\special{pa  3150  -472}\special{pa  3150    -0}%
\special{fp}%
}%
{%
\color{red}%
\special{pa  3465  -472}\special{pa  3465    -0}%
\special{fp}%
}%
{%
\color{red}%
\special{pa  3780  -472}\special{pa  3780    -0}%
\special{fp}%
}%
{%
\color{red}%
\special{pa     0  -472}\special{pa  3780  -472}%
\special{fp}%
}%
{%
\color{red}%
\special{pa     0  -236}\special{pa  3780  -236}%
\special{fp}%
}%
{%
\color{red}%
\special{pa     0    -0}\special{pa  3780    -0}%
\special{fp}%
}%
{%
\color{red}%
\settowidth{\Width}{$x$}\setlength{\Width}{-0.5\Width}%
\settoheight{\Height}{$x$}\settodepth{\Depth}{$x$}\setlength{\Height}{-0.5\Height}\setlength{\Depth}{0.5\Depth}\addtolength{\Height}{\Depth}%
\put(0.4000000,0.9000000){\hspace*{\Width}\raisebox{\Height}{$x$}}%
%
}%
{%
\color{red}%
\settowidth{\Width}{$-5$}\setlength{\Width}{-0.5\Width}%
\settoheight{\Height}{$-5$}\settodepth{\Depth}{$-5$}\setlength{\Height}{-0.5\Height}\setlength{\Depth}{0.5\Depth}\addtolength{\Height}{\Depth}%
\put(1.2000000,0.9000000){\hspace*{\Width}\raisebox{\Height}{$-5$}}%
%
}%
{%
\color{red}%
\settowidth{\Width}{$-4$}\setlength{\Width}{-0.5\Width}%
\settoheight{\Height}{$-4$}\settodepth{\Depth}{$-4$}\setlength{\Height}{-0.5\Height}\setlength{\Depth}{0.5\Depth}\addtolength{\Height}{\Depth}%
\put(2.0000000,0.9000000){\hspace*{\Width}\raisebox{\Height}{$-4$}}%
%
}%
{%
\color{red}%
\settowidth{\Width}{$-3$}\setlength{\Width}{-0.5\Width}%
\settoheight{\Height}{$-3$}\settodepth{\Depth}{$-3$}\setlength{\Height}{-0.5\Height}\setlength{\Depth}{0.5\Depth}\addtolength{\Height}{\Depth}%
\put(2.8000000,0.9000000){\hspace*{\Width}\raisebox{\Height}{$-3$}}%
%
}%
{%
\color{red}%
\settowidth{\Width}{$-2$}\setlength{\Width}{-0.5\Width}%
\settoheight{\Height}{$-2$}\settodepth{\Depth}{$-2$}\setlength{\Height}{-0.5\Height}\setlength{\Depth}{0.5\Depth}\addtolength{\Height}{\Depth}%
\put(3.6000000,0.9000000){\hspace*{\Width}\raisebox{\Height}{$-2$}}%
%
}%
{%
\color{red}%
\settowidth{\Width}{$-1$}\setlength{\Width}{-0.5\Width}%
\settoheight{\Height}{$-1$}\settodepth{\Depth}{$-1$}\setlength{\Height}{-0.5\Height}\setlength{\Depth}{0.5\Depth}\addtolength{\Height}{\Depth}%
\put(4.4000000,0.9000000){\hspace*{\Width}\raisebox{\Height}{$-1$}}%
%
}%
{%
\color{red}%
\settowidth{\Width}{$0$}\setlength{\Width}{-0.5\Width}%
\settoheight{\Height}{$0$}\settodepth{\Depth}{$0$}\setlength{\Height}{-0.5\Height}\setlength{\Depth}{0.5\Depth}\addtolength{\Height}{\Depth}%
\put(5.2000000,0.9000000){\hspace*{\Width}\raisebox{\Height}{$0$}}%
%
}%
{%
\color{red}%
\settowidth{\Width}{$1$}\setlength{\Width}{-0.5\Width}%
\settoheight{\Height}{$1$}\settodepth{\Depth}{$1$}\setlength{\Height}{-0.5\Height}\setlength{\Depth}{0.5\Depth}\addtolength{\Height}{\Depth}%
\put(6.0000000,0.9000000){\hspace*{\Width}\raisebox{\Height}{$1$}}%
%
}%
{%
\color{red}%
\settowidth{\Width}{$2$}\setlength{\Width}{-0.5\Width}%
\settoheight{\Height}{$2$}\settodepth{\Depth}{$2$}\setlength{\Height}{-0.5\Height}\setlength{\Depth}{0.5\Depth}\addtolength{\Height}{\Depth}%
\put(6.8000000,0.9000000){\hspace*{\Width}\raisebox{\Height}{$2$}}%
%
}%
{%
\color{red}%
\settowidth{\Width}{$3$}\setlength{\Width}{-0.5\Width}%
\settoheight{\Height}{$3$}\settodepth{\Depth}{$3$}\setlength{\Height}{-0.5\Height}\setlength{\Depth}{0.5\Depth}\addtolength{\Height}{\Depth}%
\put(7.6000000,0.9000000){\hspace*{\Width}\raisebox{\Height}{$3$}}%
%
}%
{%
\color{red}%
\settowidth{\Width}{$4$}\setlength{\Width}{-0.5\Width}%
\settoheight{\Height}{$4$}\settodepth{\Depth}{$4$}\setlength{\Height}{-0.5\Height}\setlength{\Depth}{0.5\Depth}\addtolength{\Height}{\Depth}%
\put(8.4000000,0.9000000){\hspace*{\Width}\raisebox{\Height}{$4$}}%
%
}%
{%
\color{red}%
\settowidth{\Width}{$5$}\setlength{\Width}{-0.5\Width}%
\settoheight{\Height}{$5$}\settodepth{\Depth}{$5$}\setlength{\Height}{-0.5\Height}\setlength{\Depth}{0.5\Depth}\addtolength{\Height}{\Depth}%
\put(9.2000000,0.9000000){\hspace*{\Width}\raisebox{\Height}{$5$}}%
%
}%
{%
\color{red}%
\settowidth{\Width}{$y$}\setlength{\Width}{-0.5\Width}%
\settoheight{\Height}{$y$}\settodepth{\Depth}{$y$}\setlength{\Height}{-0.5\Height}\setlength{\Depth}{0.5\Depth}\addtolength{\Height}{\Depth}%
\put(0.4000000,0.3000000){\hspace*{\Width}\raisebox{\Height}{$y$}}%
%
}%
{%
\color{red}%
\settowidth{\Width}{$-9$}\setlength{\Width}{-0.5\Width}%
\settoheight{\Height}{$-9$}\settodepth{\Depth}{$-9$}\setlength{\Height}{-0.5\Height}\setlength{\Depth}{0.5\Depth}\addtolength{\Height}{\Depth}%
\put(1.2000000,0.3000000){\hspace*{\Width}\raisebox{\Height}{$-9$}}%
%
}%
{%
\color{red}%
\settowidth{\Width}{$-7$}\setlength{\Width}{-0.5\Width}%
\settoheight{\Height}{$-7$}\settodepth{\Depth}{$-7$}\setlength{\Height}{-0.5\Height}\setlength{\Depth}{0.5\Depth}\addtolength{\Height}{\Depth}%
\put(2.0000000,0.3000000){\hspace*{\Width}\raisebox{\Height}{$-7$}}%
%
}%
{%
\color{red}%
\settowidth{\Width}{$-5$}\setlength{\Width}{-0.5\Width}%
\settoheight{\Height}{$-5$}\settodepth{\Depth}{$-5$}\setlength{\Height}{-0.5\Height}\setlength{\Depth}{0.5\Depth}\addtolength{\Height}{\Depth}%
\put(2.8000000,0.3000000){\hspace*{\Width}\raisebox{\Height}{$-5$}}%
%
}%
{%
\color{red}%
\settowidth{\Width}{$-3$}\setlength{\Width}{-0.5\Width}%
\settoheight{\Height}{$-3$}\settodepth{\Depth}{$-3$}\setlength{\Height}{-0.5\Height}\setlength{\Depth}{0.5\Depth}\addtolength{\Height}{\Depth}%
\put(3.6000000,0.3000000){\hspace*{\Width}\raisebox{\Height}{$-3$}}%
%
}%
{%
\color{red}%
\settowidth{\Width}{$-1$}\setlength{\Width}{-0.5\Width}%
\settoheight{\Height}{$-1$}\settodepth{\Depth}{$-1$}\setlength{\Height}{-0.5\Height}\setlength{\Depth}{0.5\Depth}\addtolength{\Height}{\Depth}%
\put(4.4000000,0.3000000){\hspace*{\Width}\raisebox{\Height}{$-1$}}%
%
}%
{%
\color{red}%
\settowidth{\Width}{$1$}\setlength{\Width}{-0.5\Width}%
\settoheight{\Height}{$1$}\settodepth{\Depth}{$1$}\setlength{\Height}{-0.5\Height}\setlength{\Depth}{0.5\Depth}\addtolength{\Height}{\Depth}%
\put(5.2000000,0.3000000){\hspace*{\Width}\raisebox{\Height}{$1$}}%
%
}%
{%
\color{red}%
\settowidth{\Width}{$3$}\setlength{\Width}{-0.5\Width}%
\settoheight{\Height}{$3$}\settodepth{\Depth}{$3$}\setlength{\Height}{-0.5\Height}\setlength{\Depth}{0.5\Depth}\addtolength{\Height}{\Depth}%
\put(6.0000000,0.3000000){\hspace*{\Width}\raisebox{\Height}{$3$}}%
%
}%
{%
\color{red}%
\settowidth{\Width}{$5$}\setlength{\Width}{-0.5\Width}%
\settoheight{\Height}{$5$}\settodepth{\Depth}{$5$}\setlength{\Height}{-0.5\Height}\setlength{\Depth}{0.5\Depth}\addtolength{\Height}{\Depth}%
\put(6.8000000,0.3000000){\hspace*{\Width}\raisebox{\Height}{$5$}}%
%
}%
{%
\color{red}%
\settowidth{\Width}{$7$}\setlength{\Width}{-0.5\Width}%
\settoheight{\Height}{$7$}\settodepth{\Depth}{$7$}\setlength{\Height}{-0.5\Height}\setlength{\Depth}{0.5\Depth}\addtolength{\Height}{\Depth}%
\put(7.6000000,0.3000000){\hspace*{\Width}\raisebox{\Height}{$7$}}%
%
}%
{%
\color{red}%
\settowidth{\Width}{$9$}\setlength{\Width}{-0.5\Width}%
\settoheight{\Height}{$9$}\settodepth{\Depth}{$9$}\setlength{\Height}{-0.5\Height}\setlength{\Depth}{0.5\Depth}\addtolength{\Height}{\Depth}%
\put(8.4000000,0.3000000){\hspace*{\Width}\raisebox{\Height}{$9$}}%
%
}%
{%
\color{red}%
\settowidth{\Width}{$11$}\setlength{\Width}{-0.5\Width}%
\settoheight{\Height}{$11$}\settodepth{\Depth}{$11$}\setlength{\Height}{-0.5\Height}\setlength{\Depth}{0.5\Depth}\addtolength{\Height}{\Depth}%
\put(9.2000000,0.3000000){\hspace*{\Width}\raisebox{\Height}{$11$}}%
%
}%
\end{picture}}%}}
\putnotes{60}{6}{\scalebox{0.5}{%%% /polytech.git/n101/fig/graphpaper3.tex 
%%% Generator=presen0601.cdy 
{\unitlength=1cm%
\begin{picture}%
(12.4,12.4)(-6.2,-6.2)%
\special{pn 8}%
%
\Large\bf\boldmath%
\small%
{%
\color[rgb]{0,0,0}%
\special{pn 4}%
\special{pa -2362 -2362}\special{pa -2362  2362}%
\special{fp}%
\special{pn 8}%
}%
{%
\color[rgb]{0,0,0}%
\special{pn 4}%
\special{pa -1969 -2362}\special{pa -1969  2362}%
\special{fp}%
\special{pn 8}%
}%
{%
\color[rgb]{0,0,0}%
\special{pn 4}%
\special{pa -1575 -2362}\special{pa -1575  2362}%
\special{fp}%
\special{pn 8}%
}%
{%
\color[rgb]{0,0,0}%
\special{pn 4}%
\special{pa -1181 -2362}\special{pa -1181  2362}%
\special{fp}%
\special{pn 8}%
}%
{%
\color[rgb]{0,0,0}%
\special{pn 4}%
\special{pa  -787 -2362}\special{pa  -787  2362}%
\special{fp}%
\special{pn 8}%
}%
{%
\color[rgb]{0,0,0}%
\special{pn 4}%
\special{pa  -394 -2362}\special{pa  -394  2362}%
\special{fp}%
\special{pn 8}%
}%
{%
\color[rgb]{0,0,0}%
\special{pn 4}%
\special{pa     0 -2362}\special{pa     0  2362}%
\special{fp}%
\special{pn 8}%
}%
{%
\color[rgb]{0,0,0}%
\special{pn 4}%
\special{pa   394 -2362}\special{pa   394  2362}%
\special{fp}%
\special{pn 8}%
}%
{%
\color[rgb]{0,0,0}%
\special{pn 4}%
\special{pa   787 -2362}\special{pa   787  2362}%
\special{fp}%
\special{pn 8}%
}%
{%
\color[rgb]{0,0,0}%
\special{pn 4}%
\special{pa  1181 -2362}\special{pa  1181  2362}%
\special{fp}%
\special{pn 8}%
}%
{%
\color[rgb]{0,0,0}%
\special{pn 4}%
\special{pa  1575 -2362}\special{pa  1575  2362}%
\special{fp}%
\special{pn 8}%
}%
{%
\color[rgb]{0,0,0}%
\special{pn 4}%
\special{pa  1969 -2362}\special{pa  1969  2362}%
\special{fp}%
\special{pn 8}%
}%
{%
\color[rgb]{0,0,0}%
\special{pn 4}%
\special{pa  2362 -2362}\special{pa  2362  2362}%
\special{fp}%
\special{pn 8}%
}%
{%
\color[rgb]{0,0,0}%
\special{pn 4}%
\special{pa -2362 -2362}\special{pa  2362 -2362}%
\special{fp}%
\special{pn 8}%
}%
{%
\color[rgb]{0,0,0}%
\special{pn 4}%
\special{pa -2362 -1969}\special{pa  2362 -1969}%
\special{fp}%
\special{pn 8}%
}%
{%
\color[rgb]{0,0,0}%
\special{pn 4}%
\special{pa -2362 -1575}\special{pa  2362 -1575}%
\special{fp}%
\special{pn 8}%
}%
{%
\color[rgb]{0,0,0}%
\special{pn 4}%
\special{pa -2362 -1181}\special{pa  2362 -1181}%
\special{fp}%
\special{pn 8}%
}%
{%
\color[rgb]{0,0,0}%
\special{pn 4}%
\special{pa -2362  -787}\special{pa  2362  -787}%
\special{fp}%
\special{pn 8}%
}%
{%
\color[rgb]{0,0,0}%
\special{pn 4}%
\special{pa -2362  -394}\special{pa  2362  -394}%
\special{fp}%
\special{pn 8}%
}%
{%
\color[rgb]{0,0,0}%
\special{pn 4}%
\special{pa -2362    -0}\special{pa  2362    -0}%
\special{fp}%
\special{pn 8}%
}%
{%
\color[rgb]{0,0,0}%
\special{pn 4}%
\special{pa -2362   394}\special{pa  2362   394}%
\special{fp}%
\special{pn 8}%
}%
{%
\color[rgb]{0,0,0}%
\special{pn 4}%
\special{pa -2362   787}\special{pa  2362   787}%
\special{fp}%
\special{pn 8}%
}%
{%
\color[rgb]{0,0,0}%
\special{pn 4}%
\special{pa -2362  1181}\special{pa  2362  1181}%
\special{fp}%
\special{pn 8}%
}%
{%
\color[rgb]{0,0,0}%
\special{pn 4}%
\special{pa -2362  1575}\special{pa  2362  1575}%
\special{fp}%
\special{pn 8}%
}%
{%
\color[rgb]{0,0,0}%
\special{pn 4}%
\special{pa -2362  1969}\special{pa  2362  1969}%
\special{fp}%
\special{pn 8}%
}%
{%
\color[rgb]{0,0,0}%
\special{pn 4}%
\special{pa -2362  2362}\special{pa  2362  2362}%
\special{fp}%
\special{pn 8}%
}%
\special{pn 4}%
{%
\color[rgb]{0,0,0}%
\special{pa -2165 2362}\special{pa -2165 2323}\special{fp}\special{pa -2165 2284}\special{pa -2165 2245}\special{fp}%
\special{pa -2165 2206}\special{pa -2165 2167}\special{fp}\special{pa -2165 2128}\special{pa -2165 2089}\special{fp}%
\special{pa -2165 2050}\special{pa -2165 2011}\special{fp}\special{pa -2165 1972}\special{pa -2165 1933}\special{fp}%
\special{pa -2165 1894}\special{pa -2165 1855}\special{fp}\special{pa -2165 1816}\special{pa -2165 1777}\special{fp}%
\special{pa -2165 1737}\special{pa -2165 1698}\special{fp}\special{pa -2165 1659}\special{pa -2165 1620}\special{fp}%
\special{pa -2165 1581}\special{pa -2165 1542}\special{fp}\special{pa -2165 1503}\special{pa -2165 1464}\special{fp}%
\special{pa -2165 1425}\special{pa -2165 1386}\special{fp}\special{pa -2165 1347}\special{pa -2165 1308}\special{fp}%
\special{pa -2165 1269}\special{pa -2165 1230}\special{fp}\special{pa -2165 1191}\special{pa -2165 1152}\special{fp}%
\special{pa -2165 1113}\special{pa -2165 1074}\special{fp}\special{pa -2165 1035}\special{pa -2165 996}\special{fp}%
\special{pa -2165 957}\special{pa -2165 918}\special{fp}\special{pa -2165 879}\special{pa -2165 839}\special{fp}%
\special{pa -2165 800}\special{pa -2165 761}\special{fp}\special{pa -2165 722}\special{pa -2165 683}\special{fp}%
\special{pa -2165 644}\special{pa -2165 605}\special{fp}\special{pa -2165 566}\special{pa -2165 527}\special{fp}%
\special{pa -2165 488}\special{pa -2165 449}\special{fp}\special{pa -2165 410}\special{pa -2165 371}\special{fp}%
\special{pa -2165 332}\special{pa -2165 293}\special{fp}\special{pa -2165 254}\special{pa -2165 215}\special{fp}%
\special{pa -2165 176}\special{pa -2165 137}\special{fp}\special{pa -2165 98}\special{pa -2165 59}\special{fp}%
\special{pa -2165 20}\special{pa -2165 -20}\special{fp}\special{pa -2165 -59}\special{pa -2165 -98}\special{fp}%
\special{pa -2165 -137}\special{pa -2165 -176}\special{fp}\special{pa -2165 -215}\special{pa -2165 -254}\special{fp}%
\special{pa -2165 -293}\special{pa -2165 -332}\special{fp}\special{pa -2165 -371}\special{pa -2165 -410}\special{fp}%
\special{pa -2165 -449}\special{pa -2165 -488}\special{fp}\special{pa -2165 -527}\special{pa -2165 -566}\special{fp}%
\special{pa -2165 -605}\special{pa -2165 -644}\special{fp}\special{pa -2165 -683}\special{pa -2165 -722}\special{fp}%
\special{pa -2165 -761}\special{pa -2165 -800}\special{fp}\special{pa -2165 -839}\special{pa -2165 -879}\special{fp}%
\special{pa -2165 -918}\special{pa -2165 -957}\special{fp}\special{pa -2165 -996}\special{pa -2165 -1035}\special{fp}%
\special{pa -2165 -1074}\special{pa -2165 -1113}\special{fp}\special{pa -2165 -1152}\special{pa -2165 -1191}\special{fp}%
\special{pa -2165 -1230}\special{pa -2165 -1269}\special{fp}\special{pa -2165 -1308}\special{pa -2165 -1347}\special{fp}%
\special{pa -2165 -1386}\special{pa -2165 -1425}\special{fp}\special{pa -2165 -1464}\special{pa -2165 -1503}\special{fp}%
\special{pa -2165 -1542}\special{pa -2165 -1581}\special{fp}\special{pa -2165 -1620}\special{pa -2165 -1659}\special{fp}%
\special{pa -2165 -1698}\special{pa -2165 -1737}\special{fp}\special{pa -2165 -1777}\special{pa -2165 -1816}\special{fp}%
\special{pa -2165 -1855}\special{pa -2165 -1894}\special{fp}\special{pa -2165 -1933}\special{pa -2165 -1972}\special{fp}%
\special{pa -2165 -2011}\special{pa -2165 -2050}\special{fp}\special{pa -2165 -2089}\special{pa -2165 -2128}\special{fp}%
\special{pa -2165 -2167}\special{pa -2165 -2206}\special{fp}\special{pa -2165 -2245}\special{pa -2165 -2284}\special{fp}%
\special{pa -2165 -2323}\special{pa -2165 -2362}\special{fp}%
%
}%
{%
\color[rgb]{0,0,0}%
\special{pa -2362 2165}\special{pa -2323 2165}\special{fp}\special{pa -2284 2165}\special{pa -2245 2165}\special{fp}%
\special{pa -2206 2165}\special{pa -2167 2165}\special{fp}\special{pa -2128 2165}\special{pa -2089 2165}\special{fp}%
\special{pa -2050 2165}\special{pa -2011 2165}\special{fp}\special{pa -1972 2165}\special{pa -1933 2165}\special{fp}%
\special{pa -1894 2165}\special{pa -1855 2165}\special{fp}\special{pa -1816 2165}\special{pa -1777 2165}\special{fp}%
\special{pa -1737 2165}\special{pa -1698 2165}\special{fp}\special{pa -1659 2165}\special{pa -1620 2165}\special{fp}%
\special{pa -1581 2165}\special{pa -1542 2165}\special{fp}\special{pa -1503 2165}\special{pa -1464 2165}\special{fp}%
\special{pa -1425 2165}\special{pa -1386 2165}\special{fp}\special{pa -1347 2165}\special{pa -1308 2165}\special{fp}%
\special{pa -1269 2165}\special{pa -1230 2165}\special{fp}\special{pa -1191 2165}\special{pa -1152 2165}\special{fp}%
\special{pa -1113 2165}\special{pa -1074 2165}\special{fp}\special{pa -1035 2165}\special{pa -996 2165}\special{fp}%
\special{pa -957 2165}\special{pa -918 2165}\special{fp}\special{pa -879 2165}\special{pa -839 2165}\special{fp}%
\special{pa -800 2165}\special{pa -761 2165}\special{fp}\special{pa -722 2165}\special{pa -683 2165}\special{fp}%
\special{pa -644 2165}\special{pa -605 2165}\special{fp}\special{pa -566 2165}\special{pa -527 2165}\special{fp}%
\special{pa -488 2165}\special{pa -449 2165}\special{fp}\special{pa -410 2165}\special{pa -371 2165}\special{fp}%
\special{pa -332 2165}\special{pa -293 2165}\special{fp}\special{pa -254 2165}\special{pa -215 2165}\special{fp}%
\special{pa -176 2165}\special{pa -137 2165}\special{fp}\special{pa -98 2165}\special{pa -59 2165}\special{fp}%
\special{pa -20 2165}\special{pa 20 2165}\special{fp}\special{pa 59 2165}\special{pa 98 2165}\special{fp}%
\special{pa 137 2165}\special{pa 176 2165}\special{fp}\special{pa 215 2165}\special{pa 254 2165}\special{fp}%
\special{pa 293 2165}\special{pa 332 2165}\special{fp}\special{pa 371 2165}\special{pa 410 2165}\special{fp}%
\special{pa 449 2165}\special{pa 488 2165}\special{fp}\special{pa 527 2165}\special{pa 566 2165}\special{fp}%
\special{pa 605 2165}\special{pa 644 2165}\special{fp}\special{pa 683 2165}\special{pa 722 2165}\special{fp}%
\special{pa 761 2165}\special{pa 800 2165}\special{fp}\special{pa 839 2165}\special{pa 879 2165}\special{fp}%
\special{pa 918 2165}\special{pa 957 2165}\special{fp}\special{pa 996 2165}\special{pa 1035 2165}\special{fp}%
\special{pa 1074 2165}\special{pa 1113 2165}\special{fp}\special{pa 1152 2165}\special{pa 1191 2165}\special{fp}%
\special{pa 1230 2165}\special{pa 1269 2165}\special{fp}\special{pa 1308 2165}\special{pa 1347 2165}\special{fp}%
\special{pa 1386 2165}\special{pa 1425 2165}\special{fp}\special{pa 1464 2165}\special{pa 1503 2165}\special{fp}%
\special{pa 1542 2165}\special{pa 1581 2165}\special{fp}\special{pa 1620 2165}\special{pa 1659 2165}\special{fp}%
\special{pa 1698 2165}\special{pa 1737 2165}\special{fp}\special{pa 1777 2165}\special{pa 1816 2165}\special{fp}%
\special{pa 1855 2165}\special{pa 1894 2165}\special{fp}\special{pa 1933 2165}\special{pa 1972 2165}\special{fp}%
\special{pa 2011 2165}\special{pa 2050 2165}\special{fp}\special{pa 2089 2165}\special{pa 2128 2165}\special{fp}%
\special{pa 2167 2165}\special{pa 2206 2165}\special{fp}\special{pa 2245 2165}\special{pa 2284 2165}\special{fp}%
\special{pa 2323 2165}\special{pa 2362 2165}\special{fp}%
%
}%
{%
\color[rgb]{0,0,0}%
\special{pa -1772 2362}\special{pa -1772 2323}\special{fp}\special{pa -1772 2284}\special{pa -1772 2245}\special{fp}%
\special{pa -1772 2206}\special{pa -1772 2167}\special{fp}\special{pa -1772 2128}\special{pa -1772 2089}\special{fp}%
\special{pa -1772 2050}\special{pa -1772 2011}\special{fp}\special{pa -1772 1972}\special{pa -1772 1933}\special{fp}%
\special{pa -1772 1894}\special{pa -1772 1855}\special{fp}\special{pa -1772 1816}\special{pa -1772 1777}\special{fp}%
\special{pa -1772 1737}\special{pa -1772 1698}\special{fp}\special{pa -1772 1659}\special{pa -1772 1620}\special{fp}%
\special{pa -1772 1581}\special{pa -1772 1542}\special{fp}\special{pa -1772 1503}\special{pa -1772 1464}\special{fp}%
\special{pa -1772 1425}\special{pa -1772 1386}\special{fp}\special{pa -1772 1347}\special{pa -1772 1308}\special{fp}%
\special{pa -1772 1269}\special{pa -1772 1230}\special{fp}\special{pa -1772 1191}\special{pa -1772 1152}\special{fp}%
\special{pa -1772 1113}\special{pa -1772 1074}\special{fp}\special{pa -1772 1035}\special{pa -1772 996}\special{fp}%
\special{pa -1772 957}\special{pa -1772 918}\special{fp}\special{pa -1772 879}\special{pa -1772 839}\special{fp}%
\special{pa -1772 800}\special{pa -1772 761}\special{fp}\special{pa -1772 722}\special{pa -1772 683}\special{fp}%
\special{pa -1772 644}\special{pa -1772 605}\special{fp}\special{pa -1772 566}\special{pa -1772 527}\special{fp}%
\special{pa -1772 488}\special{pa -1772 449}\special{fp}\special{pa -1772 410}\special{pa -1772 371}\special{fp}%
\special{pa -1772 332}\special{pa -1772 293}\special{fp}\special{pa -1772 254}\special{pa -1772 215}\special{fp}%
\special{pa -1772 176}\special{pa -1772 137}\special{fp}\special{pa -1772 98}\special{pa -1772 59}\special{fp}%
\special{pa -1772 20}\special{pa -1772 -20}\special{fp}\special{pa -1772 -59}\special{pa -1772 -98}\special{fp}%
\special{pa -1772 -137}\special{pa -1772 -176}\special{fp}\special{pa -1772 -215}\special{pa -1772 -254}\special{fp}%
\special{pa -1772 -293}\special{pa -1772 -332}\special{fp}\special{pa -1772 -371}\special{pa -1772 -410}\special{fp}%
\special{pa -1772 -449}\special{pa -1772 -488}\special{fp}\special{pa -1772 -527}\special{pa -1772 -566}\special{fp}%
\special{pa -1772 -605}\special{pa -1772 -644}\special{fp}\special{pa -1772 -683}\special{pa -1772 -722}\special{fp}%
\special{pa -1772 -761}\special{pa -1772 -800}\special{fp}\special{pa -1772 -839}\special{pa -1772 -879}\special{fp}%
\special{pa -1772 -918}\special{pa -1772 -957}\special{fp}\special{pa -1772 -996}\special{pa -1772 -1035}\special{fp}%
\special{pa -1772 -1074}\special{pa -1772 -1113}\special{fp}\special{pa -1772 -1152}\special{pa -1772 -1191}\special{fp}%
\special{pa -1772 -1230}\special{pa -1772 -1269}\special{fp}\special{pa -1772 -1308}\special{pa -1772 -1347}\special{fp}%
\special{pa -1772 -1386}\special{pa -1772 -1425}\special{fp}\special{pa -1772 -1464}\special{pa -1772 -1503}\special{fp}%
\special{pa -1772 -1542}\special{pa -1772 -1581}\special{fp}\special{pa -1772 -1620}\special{pa -1772 -1659}\special{fp}%
\special{pa -1772 -1698}\special{pa -1772 -1737}\special{fp}\special{pa -1772 -1777}\special{pa -1772 -1816}\special{fp}%
\special{pa -1772 -1855}\special{pa -1772 -1894}\special{fp}\special{pa -1772 -1933}\special{pa -1772 -1972}\special{fp}%
\special{pa -1772 -2011}\special{pa -1772 -2050}\special{fp}\special{pa -1772 -2089}\special{pa -1772 -2128}\special{fp}%
\special{pa -1772 -2167}\special{pa -1772 -2206}\special{fp}\special{pa -1772 -2245}\special{pa -1772 -2284}\special{fp}%
\special{pa -1772 -2323}\special{pa -1772 -2362}\special{fp}%
%
}%
{%
\color[rgb]{0,0,0}%
\special{pa -2362 1772}\special{pa -2323 1772}\special{fp}\special{pa -2284 1772}\special{pa -2245 1772}\special{fp}%
\special{pa -2206 1772}\special{pa -2167 1772}\special{fp}\special{pa -2128 1772}\special{pa -2089 1772}\special{fp}%
\special{pa -2050 1772}\special{pa -2011 1772}\special{fp}\special{pa -1972 1772}\special{pa -1933 1772}\special{fp}%
\special{pa -1894 1772}\special{pa -1855 1772}\special{fp}\special{pa -1816 1772}\special{pa -1777 1772}\special{fp}%
\special{pa -1737 1772}\special{pa -1698 1772}\special{fp}\special{pa -1659 1772}\special{pa -1620 1772}\special{fp}%
\special{pa -1581 1772}\special{pa -1542 1772}\special{fp}\special{pa -1503 1772}\special{pa -1464 1772}\special{fp}%
\special{pa -1425 1772}\special{pa -1386 1772}\special{fp}\special{pa -1347 1772}\special{pa -1308 1772}\special{fp}%
\special{pa -1269 1772}\special{pa -1230 1772}\special{fp}\special{pa -1191 1772}\special{pa -1152 1772}\special{fp}%
\special{pa -1113 1772}\special{pa -1074 1772}\special{fp}\special{pa -1035 1772}\special{pa -996 1772}\special{fp}%
\special{pa -957 1772}\special{pa -918 1772}\special{fp}\special{pa -879 1772}\special{pa -839 1772}\special{fp}%
\special{pa -800 1772}\special{pa -761 1772}\special{fp}\special{pa -722 1772}\special{pa -683 1772}\special{fp}%
\special{pa -644 1772}\special{pa -605 1772}\special{fp}\special{pa -566 1772}\special{pa -527 1772}\special{fp}%
\special{pa -488 1772}\special{pa -449 1772}\special{fp}\special{pa -410 1772}\special{pa -371 1772}\special{fp}%
\special{pa -332 1772}\special{pa -293 1772}\special{fp}\special{pa -254 1772}\special{pa -215 1772}\special{fp}%
\special{pa -176 1772}\special{pa -137 1772}\special{fp}\special{pa -98 1772}\special{pa -59 1772}\special{fp}%
\special{pa -20 1772}\special{pa 20 1772}\special{fp}\special{pa 59 1772}\special{pa 98 1772}\special{fp}%
\special{pa 137 1772}\special{pa 176 1772}\special{fp}\special{pa 215 1772}\special{pa 254 1772}\special{fp}%
\special{pa 293 1772}\special{pa 332 1772}\special{fp}\special{pa 371 1772}\special{pa 410 1772}\special{fp}%
\special{pa 449 1772}\special{pa 488 1772}\special{fp}\special{pa 527 1772}\special{pa 566 1772}\special{fp}%
\special{pa 605 1772}\special{pa 644 1772}\special{fp}\special{pa 683 1772}\special{pa 722 1772}\special{fp}%
\special{pa 761 1772}\special{pa 800 1772}\special{fp}\special{pa 839 1772}\special{pa 879 1772}\special{fp}%
\special{pa 918 1772}\special{pa 957 1772}\special{fp}\special{pa 996 1772}\special{pa 1035 1772}\special{fp}%
\special{pa 1074 1772}\special{pa 1113 1772}\special{fp}\special{pa 1152 1772}\special{pa 1191 1772}\special{fp}%
\special{pa 1230 1772}\special{pa 1269 1772}\special{fp}\special{pa 1308 1772}\special{pa 1347 1772}\special{fp}%
\special{pa 1386 1772}\special{pa 1425 1772}\special{fp}\special{pa 1464 1772}\special{pa 1503 1772}\special{fp}%
\special{pa 1542 1772}\special{pa 1581 1772}\special{fp}\special{pa 1620 1772}\special{pa 1659 1772}\special{fp}%
\special{pa 1698 1772}\special{pa 1737 1772}\special{fp}\special{pa 1777 1772}\special{pa 1816 1772}\special{fp}%
\special{pa 1855 1772}\special{pa 1894 1772}\special{fp}\special{pa 1933 1772}\special{pa 1972 1772}\special{fp}%
\special{pa 2011 1772}\special{pa 2050 1772}\special{fp}\special{pa 2089 1772}\special{pa 2128 1772}\special{fp}%
\special{pa 2167 1772}\special{pa 2206 1772}\special{fp}\special{pa 2245 1772}\special{pa 2284 1772}\special{fp}%
\special{pa 2323 1772}\special{pa 2362 1772}\special{fp}%
%
}%
{%
\color[rgb]{0,0,0}%
\special{pa -1378 2362}\special{pa -1378 2323}\special{fp}\special{pa -1378 2284}\special{pa -1378 2245}\special{fp}%
\special{pa -1378 2206}\special{pa -1378 2167}\special{fp}\special{pa -1378 2128}\special{pa -1378 2089}\special{fp}%
\special{pa -1378 2050}\special{pa -1378 2011}\special{fp}\special{pa -1378 1972}\special{pa -1378 1933}\special{fp}%
\special{pa -1378 1894}\special{pa -1378 1855}\special{fp}\special{pa -1378 1816}\special{pa -1378 1777}\special{fp}%
\special{pa -1378 1737}\special{pa -1378 1698}\special{fp}\special{pa -1378 1659}\special{pa -1378 1620}\special{fp}%
\special{pa -1378 1581}\special{pa -1378 1542}\special{fp}\special{pa -1378 1503}\special{pa -1378 1464}\special{fp}%
\special{pa -1378 1425}\special{pa -1378 1386}\special{fp}\special{pa -1378 1347}\special{pa -1378 1308}\special{fp}%
\special{pa -1378 1269}\special{pa -1378 1230}\special{fp}\special{pa -1378 1191}\special{pa -1378 1152}\special{fp}%
\special{pa -1378 1113}\special{pa -1378 1074}\special{fp}\special{pa -1378 1035}\special{pa -1378 996}\special{fp}%
\special{pa -1378 957}\special{pa -1378 918}\special{fp}\special{pa -1378 879}\special{pa -1378 839}\special{fp}%
\special{pa -1378 800}\special{pa -1378 761}\special{fp}\special{pa -1378 722}\special{pa -1378 683}\special{fp}%
\special{pa -1378 644}\special{pa -1378 605}\special{fp}\special{pa -1378 566}\special{pa -1378 527}\special{fp}%
\special{pa -1378 488}\special{pa -1378 449}\special{fp}\special{pa -1378 410}\special{pa -1378 371}\special{fp}%
\special{pa -1378 332}\special{pa -1378 293}\special{fp}\special{pa -1378 254}\special{pa -1378 215}\special{fp}%
\special{pa -1378 176}\special{pa -1378 137}\special{fp}\special{pa -1378 98}\special{pa -1378 59}\special{fp}%
\special{pa -1378 20}\special{pa -1378 -20}\special{fp}\special{pa -1378 -59}\special{pa -1378 -98}\special{fp}%
\special{pa -1378 -137}\special{pa -1378 -176}\special{fp}\special{pa -1378 -215}\special{pa -1378 -254}\special{fp}%
\special{pa -1378 -293}\special{pa -1378 -332}\special{fp}\special{pa -1378 -371}\special{pa -1378 -410}\special{fp}%
\special{pa -1378 -449}\special{pa -1378 -488}\special{fp}\special{pa -1378 -527}\special{pa -1378 -566}\special{fp}%
\special{pa -1378 -605}\special{pa -1378 -644}\special{fp}\special{pa -1378 -683}\special{pa -1378 -722}\special{fp}%
\special{pa -1378 -761}\special{pa -1378 -800}\special{fp}\special{pa -1378 -839}\special{pa -1378 -879}\special{fp}%
\special{pa -1378 -918}\special{pa -1378 -957}\special{fp}\special{pa -1378 -996}\special{pa -1378 -1035}\special{fp}%
\special{pa -1378 -1074}\special{pa -1378 -1113}\special{fp}\special{pa -1378 -1152}\special{pa -1378 -1191}\special{fp}%
\special{pa -1378 -1230}\special{pa -1378 -1269}\special{fp}\special{pa -1378 -1308}\special{pa -1378 -1347}\special{fp}%
\special{pa -1378 -1386}\special{pa -1378 -1425}\special{fp}\special{pa -1378 -1464}\special{pa -1378 -1503}\special{fp}%
\special{pa -1378 -1542}\special{pa -1378 -1581}\special{fp}\special{pa -1378 -1620}\special{pa -1378 -1659}\special{fp}%
\special{pa -1378 -1698}\special{pa -1378 -1737}\special{fp}\special{pa -1378 -1777}\special{pa -1378 -1816}\special{fp}%
\special{pa -1378 -1855}\special{pa -1378 -1894}\special{fp}\special{pa -1378 -1933}\special{pa -1378 -1972}\special{fp}%
\special{pa -1378 -2011}\special{pa -1378 -2050}\special{fp}\special{pa -1378 -2089}\special{pa -1378 -2128}\special{fp}%
\special{pa -1378 -2167}\special{pa -1378 -2206}\special{fp}\special{pa -1378 -2245}\special{pa -1378 -2284}\special{fp}%
\special{pa -1378 -2323}\special{pa -1378 -2362}\special{fp}%
%
}%
{%
\color[rgb]{0,0,0}%
\special{pa -2362 1378}\special{pa -2323 1378}\special{fp}\special{pa -2284 1378}\special{pa -2245 1378}\special{fp}%
\special{pa -2206 1378}\special{pa -2167 1378}\special{fp}\special{pa -2128 1378}\special{pa -2089 1378}\special{fp}%
\special{pa -2050 1378}\special{pa -2011 1378}\special{fp}\special{pa -1972 1378}\special{pa -1933 1378}\special{fp}%
\special{pa -1894 1378}\special{pa -1855 1378}\special{fp}\special{pa -1816 1378}\special{pa -1777 1378}\special{fp}%
\special{pa -1737 1378}\special{pa -1698 1378}\special{fp}\special{pa -1659 1378}\special{pa -1620 1378}\special{fp}%
\special{pa -1581 1378}\special{pa -1542 1378}\special{fp}\special{pa -1503 1378}\special{pa -1464 1378}\special{fp}%
\special{pa -1425 1378}\special{pa -1386 1378}\special{fp}\special{pa -1347 1378}\special{pa -1308 1378}\special{fp}%
\special{pa -1269 1378}\special{pa -1230 1378}\special{fp}\special{pa -1191 1378}\special{pa -1152 1378}\special{fp}%
\special{pa -1113 1378}\special{pa -1074 1378}\special{fp}\special{pa -1035 1378}\special{pa -996 1378}\special{fp}%
\special{pa -957 1378}\special{pa -918 1378}\special{fp}\special{pa -879 1378}\special{pa -839 1378}\special{fp}%
\special{pa -800 1378}\special{pa -761 1378}\special{fp}\special{pa -722 1378}\special{pa -683 1378}\special{fp}%
\special{pa -644 1378}\special{pa -605 1378}\special{fp}\special{pa -566 1378}\special{pa -527 1378}\special{fp}%
\special{pa -488 1378}\special{pa -449 1378}\special{fp}\special{pa -410 1378}\special{pa -371 1378}\special{fp}%
\special{pa -332 1378}\special{pa -293 1378}\special{fp}\special{pa -254 1378}\special{pa -215 1378}\special{fp}%
\special{pa -176 1378}\special{pa -137 1378}\special{fp}\special{pa -98 1378}\special{pa -59 1378}\special{fp}%
\special{pa -20 1378}\special{pa 20 1378}\special{fp}\special{pa 59 1378}\special{pa 98 1378}\special{fp}%
\special{pa 137 1378}\special{pa 176 1378}\special{fp}\special{pa 215 1378}\special{pa 254 1378}\special{fp}%
\special{pa 293 1378}\special{pa 332 1378}\special{fp}\special{pa 371 1378}\special{pa 410 1378}\special{fp}%
\special{pa 449 1378}\special{pa 488 1378}\special{fp}\special{pa 527 1378}\special{pa 566 1378}\special{fp}%
\special{pa 605 1378}\special{pa 644 1378}\special{fp}\special{pa 683 1378}\special{pa 722 1378}\special{fp}%
\special{pa 761 1378}\special{pa 800 1378}\special{fp}\special{pa 839 1378}\special{pa 879 1378}\special{fp}%
\special{pa 918 1378}\special{pa 957 1378}\special{fp}\special{pa 996 1378}\special{pa 1035 1378}\special{fp}%
\special{pa 1074 1378}\special{pa 1113 1378}\special{fp}\special{pa 1152 1378}\special{pa 1191 1378}\special{fp}%
\special{pa 1230 1378}\special{pa 1269 1378}\special{fp}\special{pa 1308 1378}\special{pa 1347 1378}\special{fp}%
\special{pa 1386 1378}\special{pa 1425 1378}\special{fp}\special{pa 1464 1378}\special{pa 1503 1378}\special{fp}%
\special{pa 1542 1378}\special{pa 1581 1378}\special{fp}\special{pa 1620 1378}\special{pa 1659 1378}\special{fp}%
\special{pa 1698 1378}\special{pa 1737 1378}\special{fp}\special{pa 1777 1378}\special{pa 1816 1378}\special{fp}%
\special{pa 1855 1378}\special{pa 1894 1378}\special{fp}\special{pa 1933 1378}\special{pa 1972 1378}\special{fp}%
\special{pa 2011 1378}\special{pa 2050 1378}\special{fp}\special{pa 2089 1378}\special{pa 2128 1378}\special{fp}%
\special{pa 2167 1378}\special{pa 2206 1378}\special{fp}\special{pa 2245 1378}\special{pa 2284 1378}\special{fp}%
\special{pa 2323 1378}\special{pa 2362 1378}\special{fp}%
%
}%
{%
\color[rgb]{0,0,0}%
\special{pa -984 2362}\special{pa -984 2323}\special{fp}\special{pa -984 2284}\special{pa -984 2245}\special{fp}%
\special{pa -984 2206}\special{pa -984 2167}\special{fp}\special{pa -984 2128}\special{pa -984 2089}\special{fp}%
\special{pa -984 2050}\special{pa -984 2011}\special{fp}\special{pa -984 1972}\special{pa -984 1933}\special{fp}%
\special{pa -984 1894}\special{pa -984 1855}\special{fp}\special{pa -984 1816}\special{pa -984 1777}\special{fp}%
\special{pa -984 1737}\special{pa -984 1698}\special{fp}\special{pa -984 1659}\special{pa -984 1620}\special{fp}%
\special{pa -984 1581}\special{pa -984 1542}\special{fp}\special{pa -984 1503}\special{pa -984 1464}\special{fp}%
\special{pa -984 1425}\special{pa -984 1386}\special{fp}\special{pa -984 1347}\special{pa -984 1308}\special{fp}%
\special{pa -984 1269}\special{pa -984 1230}\special{fp}\special{pa -984 1191}\special{pa -984 1152}\special{fp}%
\special{pa -984 1113}\special{pa -984 1074}\special{fp}\special{pa -984 1035}\special{pa -984 996}\special{fp}%
\special{pa -984 957}\special{pa -984 918}\special{fp}\special{pa -984 879}\special{pa -984 839}\special{fp}%
\special{pa -984 800}\special{pa -984 761}\special{fp}\special{pa -984 722}\special{pa -984 683}\special{fp}%
\special{pa -984 644}\special{pa -984 605}\special{fp}\special{pa -984 566}\special{pa -984 527}\special{fp}%
\special{pa -984 488}\special{pa -984 449}\special{fp}\special{pa -984 410}\special{pa -984 371}\special{fp}%
\special{pa -984 332}\special{pa -984 293}\special{fp}\special{pa -984 254}\special{pa -984 215}\special{fp}%
\special{pa -984 176}\special{pa -984 137}\special{fp}\special{pa -984 98}\special{pa -984 59}\special{fp}%
\special{pa -984 20}\special{pa -984 -20}\special{fp}\special{pa -984 -59}\special{pa -984 -98}\special{fp}%
\special{pa -984 -137}\special{pa -984 -176}\special{fp}\special{pa -984 -215}\special{pa -984 -254}\special{fp}%
\special{pa -984 -293}\special{pa -984 -332}\special{fp}\special{pa -984 -371}\special{pa -984 -410}\special{fp}%
\special{pa -984 -449}\special{pa -984 -488}\special{fp}\special{pa -984 -527}\special{pa -984 -566}\special{fp}%
\special{pa -984 -605}\special{pa -984 -644}\special{fp}\special{pa -984 -683}\special{pa -984 -722}\special{fp}%
\special{pa -984 -761}\special{pa -984 -800}\special{fp}\special{pa -984 -839}\special{pa -984 -879}\special{fp}%
\special{pa -984 -918}\special{pa -984 -957}\special{fp}\special{pa -984 -996}\special{pa -984 -1035}\special{fp}%
\special{pa -984 -1074}\special{pa -984 -1113}\special{fp}\special{pa -984 -1152}\special{pa -984 -1191}\special{fp}%
\special{pa -984 -1230}\special{pa -984 -1269}\special{fp}\special{pa -984 -1308}\special{pa -984 -1347}\special{fp}%
\special{pa -984 -1386}\special{pa -984 -1425}\special{fp}\special{pa -984 -1464}\special{pa -984 -1503}\special{fp}%
\special{pa -984 -1542}\special{pa -984 -1581}\special{fp}\special{pa -984 -1620}\special{pa -984 -1659}\special{fp}%
\special{pa -984 -1698}\special{pa -984 -1737}\special{fp}\special{pa -984 -1777}\special{pa -984 -1816}\special{fp}%
\special{pa -984 -1855}\special{pa -984 -1894}\special{fp}\special{pa -984 -1933}\special{pa -984 -1972}\special{fp}%
\special{pa -984 -2011}\special{pa -984 -2050}\special{fp}\special{pa -984 -2089}\special{pa -984 -2128}\special{fp}%
\special{pa -984 -2167}\special{pa -984 -2206}\special{fp}\special{pa -984 -2245}\special{pa -984 -2284}\special{fp}%
\special{pa -984 -2323}\special{pa -984 -2362}\special{fp}%
%
}%
{%
\color[rgb]{0,0,0}%
\special{pa -2362 984}\special{pa -2323 984}\special{fp}\special{pa -2284 984}\special{pa -2245 984}\special{fp}%
\special{pa -2206 984}\special{pa -2167 984}\special{fp}\special{pa -2128 984}\special{pa -2089 984}\special{fp}%
\special{pa -2050 984}\special{pa -2011 984}\special{fp}\special{pa -1972 984}\special{pa -1933 984}\special{fp}%
\special{pa -1894 984}\special{pa -1855 984}\special{fp}\special{pa -1816 984}\special{pa -1777 984}\special{fp}%
\special{pa -1737 984}\special{pa -1698 984}\special{fp}\special{pa -1659 984}\special{pa -1620 984}\special{fp}%
\special{pa -1581 984}\special{pa -1542 984}\special{fp}\special{pa -1503 984}\special{pa -1464 984}\special{fp}%
\special{pa -1425 984}\special{pa -1386 984}\special{fp}\special{pa -1347 984}\special{pa -1308 984}\special{fp}%
\special{pa -1269 984}\special{pa -1230 984}\special{fp}\special{pa -1191 984}\special{pa -1152 984}\special{fp}%
\special{pa -1113 984}\special{pa -1074 984}\special{fp}\special{pa -1035 984}\special{pa -996 984}\special{fp}%
\special{pa -957 984}\special{pa -918 984}\special{fp}\special{pa -879 984}\special{pa -839 984}\special{fp}%
\special{pa -800 984}\special{pa -761 984}\special{fp}\special{pa -722 984}\special{pa -683 984}\special{fp}%
\special{pa -644 984}\special{pa -605 984}\special{fp}\special{pa -566 984}\special{pa -527 984}\special{fp}%
\special{pa -488 984}\special{pa -449 984}\special{fp}\special{pa -410 984}\special{pa -371 984}\special{fp}%
\special{pa -332 984}\special{pa -293 984}\special{fp}\special{pa -254 984}\special{pa -215 984}\special{fp}%
\special{pa -176 984}\special{pa -137 984}\special{fp}\special{pa -98 984}\special{pa -59 984}\special{fp}%
\special{pa -20 984}\special{pa 20 984}\special{fp}\special{pa 59 984}\special{pa 98 984}\special{fp}%
\special{pa 137 984}\special{pa 176 984}\special{fp}\special{pa 215 984}\special{pa 254 984}\special{fp}%
\special{pa 293 984}\special{pa 332 984}\special{fp}\special{pa 371 984}\special{pa 410 984}\special{fp}%
\special{pa 449 984}\special{pa 488 984}\special{fp}\special{pa 527 984}\special{pa 566 984}\special{fp}%
\special{pa 605 984}\special{pa 644 984}\special{fp}\special{pa 683 984}\special{pa 722 984}\special{fp}%
\special{pa 761 984}\special{pa 800 984}\special{fp}\special{pa 839 984}\special{pa 879 984}\special{fp}%
\special{pa 918 984}\special{pa 957 984}\special{fp}\special{pa 996 984}\special{pa 1035 984}\special{fp}%
\special{pa 1074 984}\special{pa 1113 984}\special{fp}\special{pa 1152 984}\special{pa 1191 984}\special{fp}%
\special{pa 1230 984}\special{pa 1269 984}\special{fp}\special{pa 1308 984}\special{pa 1347 984}\special{fp}%
\special{pa 1386 984}\special{pa 1425 984}\special{fp}\special{pa 1464 984}\special{pa 1503 984}\special{fp}%
\special{pa 1542 984}\special{pa 1581 984}\special{fp}\special{pa 1620 984}\special{pa 1659 984}\special{fp}%
\special{pa 1698 984}\special{pa 1737 984}\special{fp}\special{pa 1777 984}\special{pa 1816 984}\special{fp}%
\special{pa 1855 984}\special{pa 1894 984}\special{fp}\special{pa 1933 984}\special{pa 1972 984}\special{fp}%
\special{pa 2011 984}\special{pa 2050 984}\special{fp}\special{pa 2089 984}\special{pa 2128 984}\special{fp}%
\special{pa 2167 984}\special{pa 2206 984}\special{fp}\special{pa 2245 984}\special{pa 2284 984}\special{fp}%
\special{pa 2323 984}\special{pa 2362 984}\special{fp}%
%
}%
{%
\color[rgb]{0,0,0}%
\special{pa -591 2362}\special{pa -591 2323}\special{fp}\special{pa -591 2284}\special{pa -591 2245}\special{fp}%
\special{pa -591 2206}\special{pa -591 2167}\special{fp}\special{pa -591 2128}\special{pa -591 2089}\special{fp}%
\special{pa -591 2050}\special{pa -591 2011}\special{fp}\special{pa -591 1972}\special{pa -591 1933}\special{fp}%
\special{pa -591 1894}\special{pa -591 1855}\special{fp}\special{pa -591 1816}\special{pa -591 1777}\special{fp}%
\special{pa -591 1737}\special{pa -591 1698}\special{fp}\special{pa -591 1659}\special{pa -591 1620}\special{fp}%
\special{pa -591 1581}\special{pa -591 1542}\special{fp}\special{pa -591 1503}\special{pa -591 1464}\special{fp}%
\special{pa -591 1425}\special{pa -591 1386}\special{fp}\special{pa -591 1347}\special{pa -591 1308}\special{fp}%
\special{pa -591 1269}\special{pa -591 1230}\special{fp}\special{pa -591 1191}\special{pa -591 1152}\special{fp}%
\special{pa -591 1113}\special{pa -591 1074}\special{fp}\special{pa -591 1035}\special{pa -591 996}\special{fp}%
\special{pa -591 957}\special{pa -591 918}\special{fp}\special{pa -591 879}\special{pa -591 839}\special{fp}%
\special{pa -591 800}\special{pa -591 761}\special{fp}\special{pa -591 722}\special{pa -591 683}\special{fp}%
\special{pa -591 644}\special{pa -591 605}\special{fp}\special{pa -591 566}\special{pa -591 527}\special{fp}%
\special{pa -591 488}\special{pa -591 449}\special{fp}\special{pa -591 410}\special{pa -591 371}\special{fp}%
\special{pa -591 332}\special{pa -591 293}\special{fp}\special{pa -591 254}\special{pa -591 215}\special{fp}%
\special{pa -591 176}\special{pa -591 137}\special{fp}\special{pa -591 98}\special{pa -591 59}\special{fp}%
\special{pa -591 20}\special{pa -591 -20}\special{fp}\special{pa -591 -59}\special{pa -591 -98}\special{fp}%
\special{pa -591 -137}\special{pa -591 -176}\special{fp}\special{pa -591 -215}\special{pa -591 -254}\special{fp}%
\special{pa -591 -293}\special{pa -591 -332}\special{fp}\special{pa -591 -371}\special{pa -591 -410}\special{fp}%
\special{pa -591 -449}\special{pa -591 -488}\special{fp}\special{pa -591 -527}\special{pa -591 -566}\special{fp}%
\special{pa -591 -605}\special{pa -591 -644}\special{fp}\special{pa -591 -683}\special{pa -591 -722}\special{fp}%
\special{pa -591 -761}\special{pa -591 -800}\special{fp}\special{pa -591 -839}\special{pa -591 -879}\special{fp}%
\special{pa -591 -918}\special{pa -591 -957}\special{fp}\special{pa -591 -996}\special{pa -591 -1035}\special{fp}%
\special{pa -591 -1074}\special{pa -591 -1113}\special{fp}\special{pa -591 -1152}\special{pa -591 -1191}\special{fp}%
\special{pa -591 -1230}\special{pa -591 -1269}\special{fp}\special{pa -591 -1308}\special{pa -591 -1347}\special{fp}%
\special{pa -591 -1386}\special{pa -591 -1425}\special{fp}\special{pa -591 -1464}\special{pa -591 -1503}\special{fp}%
\special{pa -591 -1542}\special{pa -591 -1581}\special{fp}\special{pa -591 -1620}\special{pa -591 -1659}\special{fp}%
\special{pa -591 -1698}\special{pa -591 -1737}\special{fp}\special{pa -591 -1777}\special{pa -591 -1816}\special{fp}%
\special{pa -591 -1855}\special{pa -591 -1894}\special{fp}\special{pa -591 -1933}\special{pa -591 -1972}\special{fp}%
\special{pa -591 -2011}\special{pa -591 -2050}\special{fp}\special{pa -591 -2089}\special{pa -591 -2128}\special{fp}%
\special{pa -591 -2167}\special{pa -591 -2206}\special{fp}\special{pa -591 -2245}\special{pa -591 -2284}\special{fp}%
\special{pa -591 -2323}\special{pa -591 -2362}\special{fp}%
%
}%
{%
\color[rgb]{0,0,0}%
\special{pa -2362 591}\special{pa -2323 591}\special{fp}\special{pa -2284 591}\special{pa -2245 591}\special{fp}%
\special{pa -2206 591}\special{pa -2167 591}\special{fp}\special{pa -2128 591}\special{pa -2089 591}\special{fp}%
\special{pa -2050 591}\special{pa -2011 591}\special{fp}\special{pa -1972 591}\special{pa -1933 591}\special{fp}%
\special{pa -1894 591}\special{pa -1855 591}\special{fp}\special{pa -1816 591}\special{pa -1777 591}\special{fp}%
\special{pa -1737 591}\special{pa -1698 591}\special{fp}\special{pa -1659 591}\special{pa -1620 591}\special{fp}%
\special{pa -1581 591}\special{pa -1542 591}\special{fp}\special{pa -1503 591}\special{pa -1464 591}\special{fp}%
\special{pa -1425 591}\special{pa -1386 591}\special{fp}\special{pa -1347 591}\special{pa -1308 591}\special{fp}%
\special{pa -1269 591}\special{pa -1230 591}\special{fp}\special{pa -1191 591}\special{pa -1152 591}\special{fp}%
\special{pa -1113 591}\special{pa -1074 591}\special{fp}\special{pa -1035 591}\special{pa -996 591}\special{fp}%
\special{pa -957 591}\special{pa -918 591}\special{fp}\special{pa -879 591}\special{pa -839 591}\special{fp}%
\special{pa -800 591}\special{pa -761 591}\special{fp}\special{pa -722 591}\special{pa -683 591}\special{fp}%
\special{pa -644 591}\special{pa -605 591}\special{fp}\special{pa -566 591}\special{pa -527 591}\special{fp}%
\special{pa -488 591}\special{pa -449 591}\special{fp}\special{pa -410 591}\special{pa -371 591}\special{fp}%
\special{pa -332 591}\special{pa -293 591}\special{fp}\special{pa -254 591}\special{pa -215 591}\special{fp}%
\special{pa -176 591}\special{pa -137 591}\special{fp}\special{pa -98 591}\special{pa -59 591}\special{fp}%
\special{pa -20 591}\special{pa 20 591}\special{fp}\special{pa 59 591}\special{pa 98 591}\special{fp}%
\special{pa 137 591}\special{pa 176 591}\special{fp}\special{pa 215 591}\special{pa 254 591}\special{fp}%
\special{pa 293 591}\special{pa 332 591}\special{fp}\special{pa 371 591}\special{pa 410 591}\special{fp}%
\special{pa 449 591}\special{pa 488 591}\special{fp}\special{pa 527 591}\special{pa 566 591}\special{fp}%
\special{pa 605 591}\special{pa 644 591}\special{fp}\special{pa 683 591}\special{pa 722 591}\special{fp}%
\special{pa 761 591}\special{pa 800 591}\special{fp}\special{pa 839 591}\special{pa 879 591}\special{fp}%
\special{pa 918 591}\special{pa 957 591}\special{fp}\special{pa 996 591}\special{pa 1035 591}\special{fp}%
\special{pa 1074 591}\special{pa 1113 591}\special{fp}\special{pa 1152 591}\special{pa 1191 591}\special{fp}%
\special{pa 1230 591}\special{pa 1269 591}\special{fp}\special{pa 1308 591}\special{pa 1347 591}\special{fp}%
\special{pa 1386 591}\special{pa 1425 591}\special{fp}\special{pa 1464 591}\special{pa 1503 591}\special{fp}%
\special{pa 1542 591}\special{pa 1581 591}\special{fp}\special{pa 1620 591}\special{pa 1659 591}\special{fp}%
\special{pa 1698 591}\special{pa 1737 591}\special{fp}\special{pa 1777 591}\special{pa 1816 591}\special{fp}%
\special{pa 1855 591}\special{pa 1894 591}\special{fp}\special{pa 1933 591}\special{pa 1972 591}\special{fp}%
\special{pa 2011 591}\special{pa 2050 591}\special{fp}\special{pa 2089 591}\special{pa 2128 591}\special{fp}%
\special{pa 2167 591}\special{pa 2206 591}\special{fp}\special{pa 2245 591}\special{pa 2284 591}\special{fp}%
\special{pa 2323 591}\special{pa 2362 591}\special{fp}%
%
}%
{%
\color[rgb]{0,0,0}%
\special{pa -197 2362}\special{pa -197 2323}\special{fp}\special{pa -197 2284}\special{pa -197 2245}\special{fp}%
\special{pa -197 2206}\special{pa -197 2167}\special{fp}\special{pa -197 2128}\special{pa -197 2089}\special{fp}%
\special{pa -197 2050}\special{pa -197 2011}\special{fp}\special{pa -197 1972}\special{pa -197 1933}\special{fp}%
\special{pa -197 1894}\special{pa -197 1855}\special{fp}\special{pa -197 1816}\special{pa -197 1777}\special{fp}%
\special{pa -197 1737}\special{pa -197 1698}\special{fp}\special{pa -197 1659}\special{pa -197 1620}\special{fp}%
\special{pa -197 1581}\special{pa -197 1542}\special{fp}\special{pa -197 1503}\special{pa -197 1464}\special{fp}%
\special{pa -197 1425}\special{pa -197 1386}\special{fp}\special{pa -197 1347}\special{pa -197 1308}\special{fp}%
\special{pa -197 1269}\special{pa -197 1230}\special{fp}\special{pa -197 1191}\special{pa -197 1152}\special{fp}%
\special{pa -197 1113}\special{pa -197 1074}\special{fp}\special{pa -197 1035}\special{pa -197 996}\special{fp}%
\special{pa -197 957}\special{pa -197 918}\special{fp}\special{pa -197 879}\special{pa -197 839}\special{fp}%
\special{pa -197 800}\special{pa -197 761}\special{fp}\special{pa -197 722}\special{pa -197 683}\special{fp}%
\special{pa -197 644}\special{pa -197 605}\special{fp}\special{pa -197 566}\special{pa -197 527}\special{fp}%
\special{pa -197 488}\special{pa -197 449}\special{fp}\special{pa -197 410}\special{pa -197 371}\special{fp}%
\special{pa -197 332}\special{pa -197 293}\special{fp}\special{pa -197 254}\special{pa -197 215}\special{fp}%
\special{pa -197 176}\special{pa -197 137}\special{fp}\special{pa -197 98}\special{pa -197 59}\special{fp}%
\special{pa -197 20}\special{pa -197 -20}\special{fp}\special{pa -197 -59}\special{pa -197 -98}\special{fp}%
\special{pa -197 -137}\special{pa -197 -176}\special{fp}\special{pa -197 -215}\special{pa -197 -254}\special{fp}%
\special{pa -197 -293}\special{pa -197 -332}\special{fp}\special{pa -197 -371}\special{pa -197 -410}\special{fp}%
\special{pa -197 -449}\special{pa -197 -488}\special{fp}\special{pa -197 -527}\special{pa -197 -566}\special{fp}%
\special{pa -197 -605}\special{pa -197 -644}\special{fp}\special{pa -197 -683}\special{pa -197 -722}\special{fp}%
\special{pa -197 -761}\special{pa -197 -800}\special{fp}\special{pa -197 -839}\special{pa -197 -879}\special{fp}%
\special{pa -197 -918}\special{pa -197 -957}\special{fp}\special{pa -197 -996}\special{pa -197 -1035}\special{fp}%
\special{pa -197 -1074}\special{pa -197 -1113}\special{fp}\special{pa -197 -1152}\special{pa -197 -1191}\special{fp}%
\special{pa -197 -1230}\special{pa -197 -1269}\special{fp}\special{pa -197 -1308}\special{pa -197 -1347}\special{fp}%
\special{pa -197 -1386}\special{pa -197 -1425}\special{fp}\special{pa -197 -1464}\special{pa -197 -1503}\special{fp}%
\special{pa -197 -1542}\special{pa -197 -1581}\special{fp}\special{pa -197 -1620}\special{pa -197 -1659}\special{fp}%
\special{pa -197 -1698}\special{pa -197 -1737}\special{fp}\special{pa -197 -1777}\special{pa -197 -1816}\special{fp}%
\special{pa -197 -1855}\special{pa -197 -1894}\special{fp}\special{pa -197 -1933}\special{pa -197 -1972}\special{fp}%
\special{pa -197 -2011}\special{pa -197 -2050}\special{fp}\special{pa -197 -2089}\special{pa -197 -2128}\special{fp}%
\special{pa -197 -2167}\special{pa -197 -2206}\special{fp}\special{pa -197 -2245}\special{pa -197 -2284}\special{fp}%
\special{pa -197 -2323}\special{pa -197 -2362}\special{fp}%
%
}%
{%
\color[rgb]{0,0,0}%
\special{pa -2362 197}\special{pa -2323 197}\special{fp}\special{pa -2284 197}\special{pa -2245 197}\special{fp}%
\special{pa -2206 197}\special{pa -2167 197}\special{fp}\special{pa -2128 197}\special{pa -2089 197}\special{fp}%
\special{pa -2050 197}\special{pa -2011 197}\special{fp}\special{pa -1972 197}\special{pa -1933 197}\special{fp}%
\special{pa -1894 197}\special{pa -1855 197}\special{fp}\special{pa -1816 197}\special{pa -1777 197}\special{fp}%
\special{pa -1737 197}\special{pa -1698 197}\special{fp}\special{pa -1659 197}\special{pa -1620 197}\special{fp}%
\special{pa -1581 197}\special{pa -1542 197}\special{fp}\special{pa -1503 197}\special{pa -1464 197}\special{fp}%
\special{pa -1425 197}\special{pa -1386 197}\special{fp}\special{pa -1347 197}\special{pa -1308 197}\special{fp}%
\special{pa -1269 197}\special{pa -1230 197}\special{fp}\special{pa -1191 197}\special{pa -1152 197}\special{fp}%
\special{pa -1113 197}\special{pa -1074 197}\special{fp}\special{pa -1035 197}\special{pa -996 197}\special{fp}%
\special{pa -957 197}\special{pa -918 197}\special{fp}\special{pa -879 197}\special{pa -839 197}\special{fp}%
\special{pa -800 197}\special{pa -761 197}\special{fp}\special{pa -722 197}\special{pa -683 197}\special{fp}%
\special{pa -644 197}\special{pa -605 197}\special{fp}\special{pa -566 197}\special{pa -527 197}\special{fp}%
\special{pa -488 197}\special{pa -449 197}\special{fp}\special{pa -410 197}\special{pa -371 197}\special{fp}%
\special{pa -332 197}\special{pa -293 197}\special{fp}\special{pa -254 197}\special{pa -215 197}\special{fp}%
\special{pa -176 197}\special{pa -137 197}\special{fp}\special{pa -98 197}\special{pa -59 197}\special{fp}%
\special{pa -20 197}\special{pa 20 197}\special{fp}\special{pa 59 197}\special{pa 98 197}\special{fp}%
\special{pa 137 197}\special{pa 176 197}\special{fp}\special{pa 215 197}\special{pa 254 197}\special{fp}%
\special{pa 293 197}\special{pa 332 197}\special{fp}\special{pa 371 197}\special{pa 410 197}\special{fp}%
\special{pa 449 197}\special{pa 488 197}\special{fp}\special{pa 527 197}\special{pa 566 197}\special{fp}%
\special{pa 605 197}\special{pa 644 197}\special{fp}\special{pa 683 197}\special{pa 722 197}\special{fp}%
\special{pa 761 197}\special{pa 800 197}\special{fp}\special{pa 839 197}\special{pa 879 197}\special{fp}%
\special{pa 918 197}\special{pa 957 197}\special{fp}\special{pa 996 197}\special{pa 1035 197}\special{fp}%
\special{pa 1074 197}\special{pa 1113 197}\special{fp}\special{pa 1152 197}\special{pa 1191 197}\special{fp}%
\special{pa 1230 197}\special{pa 1269 197}\special{fp}\special{pa 1308 197}\special{pa 1347 197}\special{fp}%
\special{pa 1386 197}\special{pa 1425 197}\special{fp}\special{pa 1464 197}\special{pa 1503 197}\special{fp}%
\special{pa 1542 197}\special{pa 1581 197}\special{fp}\special{pa 1620 197}\special{pa 1659 197}\special{fp}%
\special{pa 1698 197}\special{pa 1737 197}\special{fp}\special{pa 1777 197}\special{pa 1816 197}\special{fp}%
\special{pa 1855 197}\special{pa 1894 197}\special{fp}\special{pa 1933 197}\special{pa 1972 197}\special{fp}%
\special{pa 2011 197}\special{pa 2050 197}\special{fp}\special{pa 2089 197}\special{pa 2128 197}\special{fp}%
\special{pa 2167 197}\special{pa 2206 197}\special{fp}\special{pa 2245 197}\special{pa 2284 197}\special{fp}%
\special{pa 2323 197}\special{pa 2362 197}\special{fp}%
%
}%
{%
\color[rgb]{0,0,0}%
\special{pa 197 2362}\special{pa 197 2323}\special{fp}\special{pa 197 2284}\special{pa 197 2245}\special{fp}%
\special{pa 197 2206}\special{pa 197 2167}\special{fp}\special{pa 197 2128}\special{pa 197 2089}\special{fp}%
\special{pa 197 2050}\special{pa 197 2011}\special{fp}\special{pa 197 1972}\special{pa 197 1933}\special{fp}%
\special{pa 197 1894}\special{pa 197 1855}\special{fp}\special{pa 197 1816}\special{pa 197 1777}\special{fp}%
\special{pa 197 1737}\special{pa 197 1698}\special{fp}\special{pa 197 1659}\special{pa 197 1620}\special{fp}%
\special{pa 197 1581}\special{pa 197 1542}\special{fp}\special{pa 197 1503}\special{pa 197 1464}\special{fp}%
\special{pa 197 1425}\special{pa 197 1386}\special{fp}\special{pa 197 1347}\special{pa 197 1308}\special{fp}%
\special{pa 197 1269}\special{pa 197 1230}\special{fp}\special{pa 197 1191}\special{pa 197 1152}\special{fp}%
\special{pa 197 1113}\special{pa 197 1074}\special{fp}\special{pa 197 1035}\special{pa 197 996}\special{fp}%
\special{pa 197 957}\special{pa 197 918}\special{fp}\special{pa 197 879}\special{pa 197 839}\special{fp}%
\special{pa 197 800}\special{pa 197 761}\special{fp}\special{pa 197 722}\special{pa 197 683}\special{fp}%
\special{pa 197 644}\special{pa 197 605}\special{fp}\special{pa 197 566}\special{pa 197 527}\special{fp}%
\special{pa 197 488}\special{pa 197 449}\special{fp}\special{pa 197 410}\special{pa 197 371}\special{fp}%
\special{pa 197 332}\special{pa 197 293}\special{fp}\special{pa 197 254}\special{pa 197 215}\special{fp}%
\special{pa 197 176}\special{pa 197 137}\special{fp}\special{pa 197 98}\special{pa 197 59}\special{fp}%
\special{pa 197 20}\special{pa 197 -20}\special{fp}\special{pa 197 -59}\special{pa 197 -98}\special{fp}%
\special{pa 197 -137}\special{pa 197 -176}\special{fp}\special{pa 197 -215}\special{pa 197 -254}\special{fp}%
\special{pa 197 -293}\special{pa 197 -332}\special{fp}\special{pa 197 -371}\special{pa 197 -410}\special{fp}%
\special{pa 197 -449}\special{pa 197 -488}\special{fp}\special{pa 197 -527}\special{pa 197 -566}\special{fp}%
\special{pa 197 -605}\special{pa 197 -644}\special{fp}\special{pa 197 -683}\special{pa 197 -722}\special{fp}%
\special{pa 197 -761}\special{pa 197 -800}\special{fp}\special{pa 197 -839}\special{pa 197 -879}\special{fp}%
\special{pa 197 -918}\special{pa 197 -957}\special{fp}\special{pa 197 -996}\special{pa 197 -1035}\special{fp}%
\special{pa 197 -1074}\special{pa 197 -1113}\special{fp}\special{pa 197 -1152}\special{pa 197 -1191}\special{fp}%
\special{pa 197 -1230}\special{pa 197 -1269}\special{fp}\special{pa 197 -1308}\special{pa 197 -1347}\special{fp}%
\special{pa 197 -1386}\special{pa 197 -1425}\special{fp}\special{pa 197 -1464}\special{pa 197 -1503}\special{fp}%
\special{pa 197 -1542}\special{pa 197 -1581}\special{fp}\special{pa 197 -1620}\special{pa 197 -1659}\special{fp}%
\special{pa 197 -1698}\special{pa 197 -1737}\special{fp}\special{pa 197 -1777}\special{pa 197 -1816}\special{fp}%
\special{pa 197 -1855}\special{pa 197 -1894}\special{fp}\special{pa 197 -1933}\special{pa 197 -1972}\special{fp}%
\special{pa 197 -2011}\special{pa 197 -2050}\special{fp}\special{pa 197 -2089}\special{pa 197 -2128}\special{fp}%
\special{pa 197 -2167}\special{pa 197 -2206}\special{fp}\special{pa 197 -2245}\special{pa 197 -2284}\special{fp}%
\special{pa 197 -2323}\special{pa 197 -2362}\special{fp}%
%
}%
{%
\color[rgb]{0,0,0}%
\special{pa -2362 -197}\special{pa -2323 -197}\special{fp}\special{pa -2284 -197}\special{pa -2245 -197}\special{fp}%
\special{pa -2206 -197}\special{pa -2167 -197}\special{fp}\special{pa -2128 -197}\special{pa -2089 -197}\special{fp}%
\special{pa -2050 -197}\special{pa -2011 -197}\special{fp}\special{pa -1972 -197}\special{pa -1933 -197}\special{fp}%
\special{pa -1894 -197}\special{pa -1855 -197}\special{fp}\special{pa -1816 -197}\special{pa -1777 -197}\special{fp}%
\special{pa -1737 -197}\special{pa -1698 -197}\special{fp}\special{pa -1659 -197}\special{pa -1620 -197}\special{fp}%
\special{pa -1581 -197}\special{pa -1542 -197}\special{fp}\special{pa -1503 -197}\special{pa -1464 -197}\special{fp}%
\special{pa -1425 -197}\special{pa -1386 -197}\special{fp}\special{pa -1347 -197}\special{pa -1308 -197}\special{fp}%
\special{pa -1269 -197}\special{pa -1230 -197}\special{fp}\special{pa -1191 -197}\special{pa -1152 -197}\special{fp}%
\special{pa -1113 -197}\special{pa -1074 -197}\special{fp}\special{pa -1035 -197}\special{pa -996 -197}\special{fp}%
\special{pa -957 -197}\special{pa -918 -197}\special{fp}\special{pa -879 -197}\special{pa -839 -197}\special{fp}%
\special{pa -800 -197}\special{pa -761 -197}\special{fp}\special{pa -722 -197}\special{pa -683 -197}\special{fp}%
\special{pa -644 -197}\special{pa -605 -197}\special{fp}\special{pa -566 -197}\special{pa -527 -197}\special{fp}%
\special{pa -488 -197}\special{pa -449 -197}\special{fp}\special{pa -410 -197}\special{pa -371 -197}\special{fp}%
\special{pa -332 -197}\special{pa -293 -197}\special{fp}\special{pa -254 -197}\special{pa -215 -197}\special{fp}%
\special{pa -176 -197}\special{pa -137 -197}\special{fp}\special{pa -98 -197}\special{pa -59 -197}\special{fp}%
\special{pa -20 -197}\special{pa 20 -197}\special{fp}\special{pa 59 -197}\special{pa 98 -197}\special{fp}%
\special{pa 137 -197}\special{pa 176 -197}\special{fp}\special{pa 215 -197}\special{pa 254 -197}\special{fp}%
\special{pa 293 -197}\special{pa 332 -197}\special{fp}\special{pa 371 -197}\special{pa 410 -197}\special{fp}%
\special{pa 449 -197}\special{pa 488 -197}\special{fp}\special{pa 527 -197}\special{pa 566 -197}\special{fp}%
\special{pa 605 -197}\special{pa 644 -197}\special{fp}\special{pa 683 -197}\special{pa 722 -197}\special{fp}%
\special{pa 761 -197}\special{pa 800 -197}\special{fp}\special{pa 839 -197}\special{pa 879 -197}\special{fp}%
\special{pa 918 -197}\special{pa 957 -197}\special{fp}\special{pa 996 -197}\special{pa 1035 -197}\special{fp}%
\special{pa 1074 -197}\special{pa 1113 -197}\special{fp}\special{pa 1152 -197}\special{pa 1191 -197}\special{fp}%
\special{pa 1230 -197}\special{pa 1269 -197}\special{fp}\special{pa 1308 -197}\special{pa 1347 -197}\special{fp}%
\special{pa 1386 -197}\special{pa 1425 -197}\special{fp}\special{pa 1464 -197}\special{pa 1503 -197}\special{fp}%
\special{pa 1542 -197}\special{pa 1581 -197}\special{fp}\special{pa 1620 -197}\special{pa 1659 -197}\special{fp}%
\special{pa 1698 -197}\special{pa 1737 -197}\special{fp}\special{pa 1777 -197}\special{pa 1816 -197}\special{fp}%
\special{pa 1855 -197}\special{pa 1894 -197}\special{fp}\special{pa 1933 -197}\special{pa 1972 -197}\special{fp}%
\special{pa 2011 -197}\special{pa 2050 -197}\special{fp}\special{pa 2089 -197}\special{pa 2128 -197}\special{fp}%
\special{pa 2167 -197}\special{pa 2206 -197}\special{fp}\special{pa 2245 -197}\special{pa 2284 -197}\special{fp}%
\special{pa 2323 -197}\special{pa 2362 -197}\special{fp}%
%
}%
{%
\color[rgb]{0,0,0}%
\special{pa 591 2362}\special{pa 591 2323}\special{fp}\special{pa 591 2284}\special{pa 591 2245}\special{fp}%
\special{pa 591 2206}\special{pa 591 2167}\special{fp}\special{pa 591 2128}\special{pa 591 2089}\special{fp}%
\special{pa 591 2050}\special{pa 591 2011}\special{fp}\special{pa 591 1972}\special{pa 591 1933}\special{fp}%
\special{pa 591 1894}\special{pa 591 1855}\special{fp}\special{pa 591 1816}\special{pa 591 1777}\special{fp}%
\special{pa 591 1737}\special{pa 591 1698}\special{fp}\special{pa 591 1659}\special{pa 591 1620}\special{fp}%
\special{pa 591 1581}\special{pa 591 1542}\special{fp}\special{pa 591 1503}\special{pa 591 1464}\special{fp}%
\special{pa 591 1425}\special{pa 591 1386}\special{fp}\special{pa 591 1347}\special{pa 591 1308}\special{fp}%
\special{pa 591 1269}\special{pa 591 1230}\special{fp}\special{pa 591 1191}\special{pa 591 1152}\special{fp}%
\special{pa 591 1113}\special{pa 591 1074}\special{fp}\special{pa 591 1035}\special{pa 591 996}\special{fp}%
\special{pa 591 957}\special{pa 591 918}\special{fp}\special{pa 591 879}\special{pa 591 839}\special{fp}%
\special{pa 591 800}\special{pa 591 761}\special{fp}\special{pa 591 722}\special{pa 591 683}\special{fp}%
\special{pa 591 644}\special{pa 591 605}\special{fp}\special{pa 591 566}\special{pa 591 527}\special{fp}%
\special{pa 591 488}\special{pa 591 449}\special{fp}\special{pa 591 410}\special{pa 591 371}\special{fp}%
\special{pa 591 332}\special{pa 591 293}\special{fp}\special{pa 591 254}\special{pa 591 215}\special{fp}%
\special{pa 591 176}\special{pa 591 137}\special{fp}\special{pa 591 98}\special{pa 591 59}\special{fp}%
\special{pa 591 20}\special{pa 591 -20}\special{fp}\special{pa 591 -59}\special{pa 591 -98}\special{fp}%
\special{pa 591 -137}\special{pa 591 -176}\special{fp}\special{pa 591 -215}\special{pa 591 -254}\special{fp}%
\special{pa 591 -293}\special{pa 591 -332}\special{fp}\special{pa 591 -371}\special{pa 591 -410}\special{fp}%
\special{pa 591 -449}\special{pa 591 -488}\special{fp}\special{pa 591 -527}\special{pa 591 -566}\special{fp}%
\special{pa 591 -605}\special{pa 591 -644}\special{fp}\special{pa 591 -683}\special{pa 591 -722}\special{fp}%
\special{pa 591 -761}\special{pa 591 -800}\special{fp}\special{pa 591 -839}\special{pa 591 -879}\special{fp}%
\special{pa 591 -918}\special{pa 591 -957}\special{fp}\special{pa 591 -996}\special{pa 591 -1035}\special{fp}%
\special{pa 591 -1074}\special{pa 591 -1113}\special{fp}\special{pa 591 -1152}\special{pa 591 -1191}\special{fp}%
\special{pa 591 -1230}\special{pa 591 -1269}\special{fp}\special{pa 591 -1308}\special{pa 591 -1347}\special{fp}%
\special{pa 591 -1386}\special{pa 591 -1425}\special{fp}\special{pa 591 -1464}\special{pa 591 -1503}\special{fp}%
\special{pa 591 -1542}\special{pa 591 -1581}\special{fp}\special{pa 591 -1620}\special{pa 591 -1659}\special{fp}%
\special{pa 591 -1698}\special{pa 591 -1737}\special{fp}\special{pa 591 -1777}\special{pa 591 -1816}\special{fp}%
\special{pa 591 -1855}\special{pa 591 -1894}\special{fp}\special{pa 591 -1933}\special{pa 591 -1972}\special{fp}%
\special{pa 591 -2011}\special{pa 591 -2050}\special{fp}\special{pa 591 -2089}\special{pa 591 -2128}\special{fp}%
\special{pa 591 -2167}\special{pa 591 -2206}\special{fp}\special{pa 591 -2245}\special{pa 591 -2284}\special{fp}%
\special{pa 591 -2323}\special{pa 591 -2362}\special{fp}%
%
}%
{%
\color[rgb]{0,0,0}%
\special{pa -2362 -591}\special{pa -2323 -591}\special{fp}\special{pa -2284 -591}\special{pa -2245 -591}\special{fp}%
\special{pa -2206 -591}\special{pa -2167 -591}\special{fp}\special{pa -2128 -591}\special{pa -2089 -591}\special{fp}%
\special{pa -2050 -591}\special{pa -2011 -591}\special{fp}\special{pa -1972 -591}\special{pa -1933 -591}\special{fp}%
\special{pa -1894 -591}\special{pa -1855 -591}\special{fp}\special{pa -1816 -591}\special{pa -1777 -591}\special{fp}%
\special{pa -1737 -591}\special{pa -1698 -591}\special{fp}\special{pa -1659 -591}\special{pa -1620 -591}\special{fp}%
\special{pa -1581 -591}\special{pa -1542 -591}\special{fp}\special{pa -1503 -591}\special{pa -1464 -591}\special{fp}%
\special{pa -1425 -591}\special{pa -1386 -591}\special{fp}\special{pa -1347 -591}\special{pa -1308 -591}\special{fp}%
\special{pa -1269 -591}\special{pa -1230 -591}\special{fp}\special{pa -1191 -591}\special{pa -1152 -591}\special{fp}%
\special{pa -1113 -591}\special{pa -1074 -591}\special{fp}\special{pa -1035 -591}\special{pa -996 -591}\special{fp}%
\special{pa -957 -591}\special{pa -918 -591}\special{fp}\special{pa -879 -591}\special{pa -839 -591}\special{fp}%
\special{pa -800 -591}\special{pa -761 -591}\special{fp}\special{pa -722 -591}\special{pa -683 -591}\special{fp}%
\special{pa -644 -591}\special{pa -605 -591}\special{fp}\special{pa -566 -591}\special{pa -527 -591}\special{fp}%
\special{pa -488 -591}\special{pa -449 -591}\special{fp}\special{pa -410 -591}\special{pa -371 -591}\special{fp}%
\special{pa -332 -591}\special{pa -293 -591}\special{fp}\special{pa -254 -591}\special{pa -215 -591}\special{fp}%
\special{pa -176 -591}\special{pa -137 -591}\special{fp}\special{pa -98 -591}\special{pa -59 -591}\special{fp}%
\special{pa -20 -591}\special{pa 20 -591}\special{fp}\special{pa 59 -591}\special{pa 98 -591}\special{fp}%
\special{pa 137 -591}\special{pa 176 -591}\special{fp}\special{pa 215 -591}\special{pa 254 -591}\special{fp}%
\special{pa 293 -591}\special{pa 332 -591}\special{fp}\special{pa 371 -591}\special{pa 410 -591}\special{fp}%
\special{pa 449 -591}\special{pa 488 -591}\special{fp}\special{pa 527 -591}\special{pa 566 -591}\special{fp}%
\special{pa 605 -591}\special{pa 644 -591}\special{fp}\special{pa 683 -591}\special{pa 722 -591}\special{fp}%
\special{pa 761 -591}\special{pa 800 -591}\special{fp}\special{pa 839 -591}\special{pa 879 -591}\special{fp}%
\special{pa 918 -591}\special{pa 957 -591}\special{fp}\special{pa 996 -591}\special{pa 1035 -591}\special{fp}%
\special{pa 1074 -591}\special{pa 1113 -591}\special{fp}\special{pa 1152 -591}\special{pa 1191 -591}\special{fp}%
\special{pa 1230 -591}\special{pa 1269 -591}\special{fp}\special{pa 1308 -591}\special{pa 1347 -591}\special{fp}%
\special{pa 1386 -591}\special{pa 1425 -591}\special{fp}\special{pa 1464 -591}\special{pa 1503 -591}\special{fp}%
\special{pa 1542 -591}\special{pa 1581 -591}\special{fp}\special{pa 1620 -591}\special{pa 1659 -591}\special{fp}%
\special{pa 1698 -591}\special{pa 1737 -591}\special{fp}\special{pa 1777 -591}\special{pa 1816 -591}\special{fp}%
\special{pa 1855 -591}\special{pa 1894 -591}\special{fp}\special{pa 1933 -591}\special{pa 1972 -591}\special{fp}%
\special{pa 2011 -591}\special{pa 2050 -591}\special{fp}\special{pa 2089 -591}\special{pa 2128 -591}\special{fp}%
\special{pa 2167 -591}\special{pa 2206 -591}\special{fp}\special{pa 2245 -591}\special{pa 2284 -591}\special{fp}%
\special{pa 2323 -591}\special{pa 2362 -591}\special{fp}%
%
}%
{%
\color[rgb]{0,0,0}%
\special{pa 984 2362}\special{pa 984 2323}\special{fp}\special{pa 984 2284}\special{pa 984 2245}\special{fp}%
\special{pa 984 2206}\special{pa 984 2167}\special{fp}\special{pa 984 2128}\special{pa 984 2089}\special{fp}%
\special{pa 984 2050}\special{pa 984 2011}\special{fp}\special{pa 984 1972}\special{pa 984 1933}\special{fp}%
\special{pa 984 1894}\special{pa 984 1855}\special{fp}\special{pa 984 1816}\special{pa 984 1777}\special{fp}%
\special{pa 984 1737}\special{pa 984 1698}\special{fp}\special{pa 984 1659}\special{pa 984 1620}\special{fp}%
\special{pa 984 1581}\special{pa 984 1542}\special{fp}\special{pa 984 1503}\special{pa 984 1464}\special{fp}%
\special{pa 984 1425}\special{pa 984 1386}\special{fp}\special{pa 984 1347}\special{pa 984 1308}\special{fp}%
\special{pa 984 1269}\special{pa 984 1230}\special{fp}\special{pa 984 1191}\special{pa 984 1152}\special{fp}%
\special{pa 984 1113}\special{pa 984 1074}\special{fp}\special{pa 984 1035}\special{pa 984 996}\special{fp}%
\special{pa 984 957}\special{pa 984 918}\special{fp}\special{pa 984 879}\special{pa 984 839}\special{fp}%
\special{pa 984 800}\special{pa 984 761}\special{fp}\special{pa 984 722}\special{pa 984 683}\special{fp}%
\special{pa 984 644}\special{pa 984 605}\special{fp}\special{pa 984 566}\special{pa 984 527}\special{fp}%
\special{pa 984 488}\special{pa 984 449}\special{fp}\special{pa 984 410}\special{pa 984 371}\special{fp}%
\special{pa 984 332}\special{pa 984 293}\special{fp}\special{pa 984 254}\special{pa 984 215}\special{fp}%
\special{pa 984 176}\special{pa 984 137}\special{fp}\special{pa 984 98}\special{pa 984 59}\special{fp}%
\special{pa 984 20}\special{pa 984 -20}\special{fp}\special{pa 984 -59}\special{pa 984 -98}\special{fp}%
\special{pa 984 -137}\special{pa 984 -176}\special{fp}\special{pa 984 -215}\special{pa 984 -254}\special{fp}%
\special{pa 984 -293}\special{pa 984 -332}\special{fp}\special{pa 984 -371}\special{pa 984 -410}\special{fp}%
\special{pa 984 -449}\special{pa 984 -488}\special{fp}\special{pa 984 -527}\special{pa 984 -566}\special{fp}%
\special{pa 984 -605}\special{pa 984 -644}\special{fp}\special{pa 984 -683}\special{pa 984 -722}\special{fp}%
\special{pa 984 -761}\special{pa 984 -800}\special{fp}\special{pa 984 -839}\special{pa 984 -879}\special{fp}%
\special{pa 984 -918}\special{pa 984 -957}\special{fp}\special{pa 984 -996}\special{pa 984 -1035}\special{fp}%
\special{pa 984 -1074}\special{pa 984 -1113}\special{fp}\special{pa 984 -1152}\special{pa 984 -1191}\special{fp}%
\special{pa 984 -1230}\special{pa 984 -1269}\special{fp}\special{pa 984 -1308}\special{pa 984 -1347}\special{fp}%
\special{pa 984 -1386}\special{pa 984 -1425}\special{fp}\special{pa 984 -1464}\special{pa 984 -1503}\special{fp}%
\special{pa 984 -1542}\special{pa 984 -1581}\special{fp}\special{pa 984 -1620}\special{pa 984 -1659}\special{fp}%
\special{pa 984 -1698}\special{pa 984 -1737}\special{fp}\special{pa 984 -1777}\special{pa 984 -1816}\special{fp}%
\special{pa 984 -1855}\special{pa 984 -1894}\special{fp}\special{pa 984 -1933}\special{pa 984 -1972}\special{fp}%
\special{pa 984 -2011}\special{pa 984 -2050}\special{fp}\special{pa 984 -2089}\special{pa 984 -2128}\special{fp}%
\special{pa 984 -2167}\special{pa 984 -2206}\special{fp}\special{pa 984 -2245}\special{pa 984 -2284}\special{fp}%
\special{pa 984 -2323}\special{pa 984 -2362}\special{fp}%
%
}%
{%
\color[rgb]{0,0,0}%
\special{pa -2362 -984}\special{pa -2323 -984}\special{fp}\special{pa -2284 -984}\special{pa -2245 -984}\special{fp}%
\special{pa -2206 -984}\special{pa -2167 -984}\special{fp}\special{pa -2128 -984}\special{pa -2089 -984}\special{fp}%
\special{pa -2050 -984}\special{pa -2011 -984}\special{fp}\special{pa -1972 -984}\special{pa -1933 -984}\special{fp}%
\special{pa -1894 -984}\special{pa -1855 -984}\special{fp}\special{pa -1816 -984}\special{pa -1777 -984}\special{fp}%
\special{pa -1737 -984}\special{pa -1698 -984}\special{fp}\special{pa -1659 -984}\special{pa -1620 -984}\special{fp}%
\special{pa -1581 -984}\special{pa -1542 -984}\special{fp}\special{pa -1503 -984}\special{pa -1464 -984}\special{fp}%
\special{pa -1425 -984}\special{pa -1386 -984}\special{fp}\special{pa -1347 -984}\special{pa -1308 -984}\special{fp}%
\special{pa -1269 -984}\special{pa -1230 -984}\special{fp}\special{pa -1191 -984}\special{pa -1152 -984}\special{fp}%
\special{pa -1113 -984}\special{pa -1074 -984}\special{fp}\special{pa -1035 -984}\special{pa -996 -984}\special{fp}%
\special{pa -957 -984}\special{pa -918 -984}\special{fp}\special{pa -879 -984}\special{pa -839 -984}\special{fp}%
\special{pa -800 -984}\special{pa -761 -984}\special{fp}\special{pa -722 -984}\special{pa -683 -984}\special{fp}%
\special{pa -644 -984}\special{pa -605 -984}\special{fp}\special{pa -566 -984}\special{pa -527 -984}\special{fp}%
\special{pa -488 -984}\special{pa -449 -984}\special{fp}\special{pa -410 -984}\special{pa -371 -984}\special{fp}%
\special{pa -332 -984}\special{pa -293 -984}\special{fp}\special{pa -254 -984}\special{pa -215 -984}\special{fp}%
\special{pa -176 -984}\special{pa -137 -984}\special{fp}\special{pa -98 -984}\special{pa -59 -984}\special{fp}%
\special{pa -20 -984}\special{pa 20 -984}\special{fp}\special{pa 59 -984}\special{pa 98 -984}\special{fp}%
\special{pa 137 -984}\special{pa 176 -984}\special{fp}\special{pa 215 -984}\special{pa 254 -984}\special{fp}%
\special{pa 293 -984}\special{pa 332 -984}\special{fp}\special{pa 371 -984}\special{pa 410 -984}\special{fp}%
\special{pa 449 -984}\special{pa 488 -984}\special{fp}\special{pa 527 -984}\special{pa 566 -984}\special{fp}%
\special{pa 605 -984}\special{pa 644 -984}\special{fp}\special{pa 683 -984}\special{pa 722 -984}\special{fp}%
\special{pa 761 -984}\special{pa 800 -984}\special{fp}\special{pa 839 -984}\special{pa 879 -984}\special{fp}%
\special{pa 918 -984}\special{pa 957 -984}\special{fp}\special{pa 996 -984}\special{pa 1035 -984}\special{fp}%
\special{pa 1074 -984}\special{pa 1113 -984}\special{fp}\special{pa 1152 -984}\special{pa 1191 -984}\special{fp}%
\special{pa 1230 -984}\special{pa 1269 -984}\special{fp}\special{pa 1308 -984}\special{pa 1347 -984}\special{fp}%
\special{pa 1386 -984}\special{pa 1425 -984}\special{fp}\special{pa 1464 -984}\special{pa 1503 -984}\special{fp}%
\special{pa 1542 -984}\special{pa 1581 -984}\special{fp}\special{pa 1620 -984}\special{pa 1659 -984}\special{fp}%
\special{pa 1698 -984}\special{pa 1737 -984}\special{fp}\special{pa 1777 -984}\special{pa 1816 -984}\special{fp}%
\special{pa 1855 -984}\special{pa 1894 -984}\special{fp}\special{pa 1933 -984}\special{pa 1972 -984}\special{fp}%
\special{pa 2011 -984}\special{pa 2050 -984}\special{fp}\special{pa 2089 -984}\special{pa 2128 -984}\special{fp}%
\special{pa 2167 -984}\special{pa 2206 -984}\special{fp}\special{pa 2245 -984}\special{pa 2284 -984}\special{fp}%
\special{pa 2323 -984}\special{pa 2362 -984}\special{fp}%
%
}%
{%
\color[rgb]{0,0,0}%
\special{pa 1378 2362}\special{pa 1378 2323}\special{fp}\special{pa 1378 2284}\special{pa 1378 2245}\special{fp}%
\special{pa 1378 2206}\special{pa 1378 2167}\special{fp}\special{pa 1378 2128}\special{pa 1378 2089}\special{fp}%
\special{pa 1378 2050}\special{pa 1378 2011}\special{fp}\special{pa 1378 1972}\special{pa 1378 1933}\special{fp}%
\special{pa 1378 1894}\special{pa 1378 1855}\special{fp}\special{pa 1378 1816}\special{pa 1378 1777}\special{fp}%
\special{pa 1378 1737}\special{pa 1378 1698}\special{fp}\special{pa 1378 1659}\special{pa 1378 1620}\special{fp}%
\special{pa 1378 1581}\special{pa 1378 1542}\special{fp}\special{pa 1378 1503}\special{pa 1378 1464}\special{fp}%
\special{pa 1378 1425}\special{pa 1378 1386}\special{fp}\special{pa 1378 1347}\special{pa 1378 1308}\special{fp}%
\special{pa 1378 1269}\special{pa 1378 1230}\special{fp}\special{pa 1378 1191}\special{pa 1378 1152}\special{fp}%
\special{pa 1378 1113}\special{pa 1378 1074}\special{fp}\special{pa 1378 1035}\special{pa 1378 996}\special{fp}%
\special{pa 1378 957}\special{pa 1378 918}\special{fp}\special{pa 1378 879}\special{pa 1378 839}\special{fp}%
\special{pa 1378 800}\special{pa 1378 761}\special{fp}\special{pa 1378 722}\special{pa 1378 683}\special{fp}%
\special{pa 1378 644}\special{pa 1378 605}\special{fp}\special{pa 1378 566}\special{pa 1378 527}\special{fp}%
\special{pa 1378 488}\special{pa 1378 449}\special{fp}\special{pa 1378 410}\special{pa 1378 371}\special{fp}%
\special{pa 1378 332}\special{pa 1378 293}\special{fp}\special{pa 1378 254}\special{pa 1378 215}\special{fp}%
\special{pa 1378 176}\special{pa 1378 137}\special{fp}\special{pa 1378 98}\special{pa 1378 59}\special{fp}%
\special{pa 1378 20}\special{pa 1378 -20}\special{fp}\special{pa 1378 -59}\special{pa 1378 -98}\special{fp}%
\special{pa 1378 -137}\special{pa 1378 -176}\special{fp}\special{pa 1378 -215}\special{pa 1378 -254}\special{fp}%
\special{pa 1378 -293}\special{pa 1378 -332}\special{fp}\special{pa 1378 -371}\special{pa 1378 -410}\special{fp}%
\special{pa 1378 -449}\special{pa 1378 -488}\special{fp}\special{pa 1378 -527}\special{pa 1378 -566}\special{fp}%
\special{pa 1378 -605}\special{pa 1378 -644}\special{fp}\special{pa 1378 -683}\special{pa 1378 -722}\special{fp}%
\special{pa 1378 -761}\special{pa 1378 -800}\special{fp}\special{pa 1378 -839}\special{pa 1378 -879}\special{fp}%
\special{pa 1378 -918}\special{pa 1378 -957}\special{fp}\special{pa 1378 -996}\special{pa 1378 -1035}\special{fp}%
\special{pa 1378 -1074}\special{pa 1378 -1113}\special{fp}\special{pa 1378 -1152}\special{pa 1378 -1191}\special{fp}%
\special{pa 1378 -1230}\special{pa 1378 -1269}\special{fp}\special{pa 1378 -1308}\special{pa 1378 -1347}\special{fp}%
\special{pa 1378 -1386}\special{pa 1378 -1425}\special{fp}\special{pa 1378 -1464}\special{pa 1378 -1503}\special{fp}%
\special{pa 1378 -1542}\special{pa 1378 -1581}\special{fp}\special{pa 1378 -1620}\special{pa 1378 -1659}\special{fp}%
\special{pa 1378 -1698}\special{pa 1378 -1737}\special{fp}\special{pa 1378 -1777}\special{pa 1378 -1816}\special{fp}%
\special{pa 1378 -1855}\special{pa 1378 -1894}\special{fp}\special{pa 1378 -1933}\special{pa 1378 -1972}\special{fp}%
\special{pa 1378 -2011}\special{pa 1378 -2050}\special{fp}\special{pa 1378 -2089}\special{pa 1378 -2128}\special{fp}%
\special{pa 1378 -2167}\special{pa 1378 -2206}\special{fp}\special{pa 1378 -2245}\special{pa 1378 -2284}\special{fp}%
\special{pa 1378 -2323}\special{pa 1378 -2362}\special{fp}%
%
}%
{%
\color[rgb]{0,0,0}%
\special{pa -2362 -1378}\special{pa -2323 -1378}\special{fp}\special{pa -2284 -1378}\special{pa -2245 -1378}\special{fp}%
\special{pa -2206 -1378}\special{pa -2167 -1378}\special{fp}\special{pa -2128 -1378}\special{pa -2089 -1378}\special{fp}%
\special{pa -2050 -1378}\special{pa -2011 -1378}\special{fp}\special{pa -1972 -1378}\special{pa -1933 -1378}\special{fp}%
\special{pa -1894 -1378}\special{pa -1855 -1378}\special{fp}\special{pa -1816 -1378}\special{pa -1777 -1378}\special{fp}%
\special{pa -1737 -1378}\special{pa -1698 -1378}\special{fp}\special{pa -1659 -1378}\special{pa -1620 -1378}\special{fp}%
\special{pa -1581 -1378}\special{pa -1542 -1378}\special{fp}\special{pa -1503 -1378}\special{pa -1464 -1378}\special{fp}%
\special{pa -1425 -1378}\special{pa -1386 -1378}\special{fp}\special{pa -1347 -1378}\special{pa -1308 -1378}\special{fp}%
\special{pa -1269 -1378}\special{pa -1230 -1378}\special{fp}\special{pa -1191 -1378}\special{pa -1152 -1378}\special{fp}%
\special{pa -1113 -1378}\special{pa -1074 -1378}\special{fp}\special{pa -1035 -1378}\special{pa -996 -1378}\special{fp}%
\special{pa -957 -1378}\special{pa -918 -1378}\special{fp}\special{pa -879 -1378}\special{pa -839 -1378}\special{fp}%
\special{pa -800 -1378}\special{pa -761 -1378}\special{fp}\special{pa -722 -1378}\special{pa -683 -1378}\special{fp}%
\special{pa -644 -1378}\special{pa -605 -1378}\special{fp}\special{pa -566 -1378}\special{pa -527 -1378}\special{fp}%
\special{pa -488 -1378}\special{pa -449 -1378}\special{fp}\special{pa -410 -1378}\special{pa -371 -1378}\special{fp}%
\special{pa -332 -1378}\special{pa -293 -1378}\special{fp}\special{pa -254 -1378}\special{pa -215 -1378}\special{fp}%
\special{pa -176 -1378}\special{pa -137 -1378}\special{fp}\special{pa -98 -1378}\special{pa -59 -1378}\special{fp}%
\special{pa -20 -1378}\special{pa 20 -1378}\special{fp}\special{pa 59 -1378}\special{pa 98 -1378}\special{fp}%
\special{pa 137 -1378}\special{pa 176 -1378}\special{fp}\special{pa 215 -1378}\special{pa 254 -1378}\special{fp}%
\special{pa 293 -1378}\special{pa 332 -1378}\special{fp}\special{pa 371 -1378}\special{pa 410 -1378}\special{fp}%
\special{pa 449 -1378}\special{pa 488 -1378}\special{fp}\special{pa 527 -1378}\special{pa 566 -1378}\special{fp}%
\special{pa 605 -1378}\special{pa 644 -1378}\special{fp}\special{pa 683 -1378}\special{pa 722 -1378}\special{fp}%
\special{pa 761 -1378}\special{pa 800 -1378}\special{fp}\special{pa 839 -1378}\special{pa 879 -1378}\special{fp}%
\special{pa 918 -1378}\special{pa 957 -1378}\special{fp}\special{pa 996 -1378}\special{pa 1035 -1378}\special{fp}%
\special{pa 1074 -1378}\special{pa 1113 -1378}\special{fp}\special{pa 1152 -1378}\special{pa 1191 -1378}\special{fp}%
\special{pa 1230 -1378}\special{pa 1269 -1378}\special{fp}\special{pa 1308 -1378}\special{pa 1347 -1378}\special{fp}%
\special{pa 1386 -1378}\special{pa 1425 -1378}\special{fp}\special{pa 1464 -1378}\special{pa 1503 -1378}\special{fp}%
\special{pa 1542 -1378}\special{pa 1581 -1378}\special{fp}\special{pa 1620 -1378}\special{pa 1659 -1378}\special{fp}%
\special{pa 1698 -1378}\special{pa 1737 -1378}\special{fp}\special{pa 1777 -1378}\special{pa 1816 -1378}\special{fp}%
\special{pa 1855 -1378}\special{pa 1894 -1378}\special{fp}\special{pa 1933 -1378}\special{pa 1972 -1378}\special{fp}%
\special{pa 2011 -1378}\special{pa 2050 -1378}\special{fp}\special{pa 2089 -1378}\special{pa 2128 -1378}\special{fp}%
\special{pa 2167 -1378}\special{pa 2206 -1378}\special{fp}\special{pa 2245 -1378}\special{pa 2284 -1378}\special{fp}%
\special{pa 2323 -1378}\special{pa 2362 -1378}\special{fp}%
%
}%
{%
\color[rgb]{0,0,0}%
\special{pa 1772 2362}\special{pa 1772 2323}\special{fp}\special{pa 1772 2284}\special{pa 1772 2245}\special{fp}%
\special{pa 1772 2206}\special{pa 1772 2167}\special{fp}\special{pa 1772 2128}\special{pa 1772 2089}\special{fp}%
\special{pa 1772 2050}\special{pa 1772 2011}\special{fp}\special{pa 1772 1972}\special{pa 1772 1933}\special{fp}%
\special{pa 1772 1894}\special{pa 1772 1855}\special{fp}\special{pa 1772 1816}\special{pa 1772 1777}\special{fp}%
\special{pa 1772 1737}\special{pa 1772 1698}\special{fp}\special{pa 1772 1659}\special{pa 1772 1620}\special{fp}%
\special{pa 1772 1581}\special{pa 1772 1542}\special{fp}\special{pa 1772 1503}\special{pa 1772 1464}\special{fp}%
\special{pa 1772 1425}\special{pa 1772 1386}\special{fp}\special{pa 1772 1347}\special{pa 1772 1308}\special{fp}%
\special{pa 1772 1269}\special{pa 1772 1230}\special{fp}\special{pa 1772 1191}\special{pa 1772 1152}\special{fp}%
\special{pa 1772 1113}\special{pa 1772 1074}\special{fp}\special{pa 1772 1035}\special{pa 1772 996}\special{fp}%
\special{pa 1772 957}\special{pa 1772 918}\special{fp}\special{pa 1772 879}\special{pa 1772 839}\special{fp}%
\special{pa 1772 800}\special{pa 1772 761}\special{fp}\special{pa 1772 722}\special{pa 1772 683}\special{fp}%
\special{pa 1772 644}\special{pa 1772 605}\special{fp}\special{pa 1772 566}\special{pa 1772 527}\special{fp}%
\special{pa 1772 488}\special{pa 1772 449}\special{fp}\special{pa 1772 410}\special{pa 1772 371}\special{fp}%
\special{pa 1772 332}\special{pa 1772 293}\special{fp}\special{pa 1772 254}\special{pa 1772 215}\special{fp}%
\special{pa 1772 176}\special{pa 1772 137}\special{fp}\special{pa 1772 98}\special{pa 1772 59}\special{fp}%
\special{pa 1772 20}\special{pa 1772 -20}\special{fp}\special{pa 1772 -59}\special{pa 1772 -98}\special{fp}%
\special{pa 1772 -137}\special{pa 1772 -176}\special{fp}\special{pa 1772 -215}\special{pa 1772 -254}\special{fp}%
\special{pa 1772 -293}\special{pa 1772 -332}\special{fp}\special{pa 1772 -371}\special{pa 1772 -410}\special{fp}%
\special{pa 1772 -449}\special{pa 1772 -488}\special{fp}\special{pa 1772 -527}\special{pa 1772 -566}\special{fp}%
\special{pa 1772 -605}\special{pa 1772 -644}\special{fp}\special{pa 1772 -683}\special{pa 1772 -722}\special{fp}%
\special{pa 1772 -761}\special{pa 1772 -800}\special{fp}\special{pa 1772 -839}\special{pa 1772 -879}\special{fp}%
\special{pa 1772 -918}\special{pa 1772 -957}\special{fp}\special{pa 1772 -996}\special{pa 1772 -1035}\special{fp}%
\special{pa 1772 -1074}\special{pa 1772 -1113}\special{fp}\special{pa 1772 -1152}\special{pa 1772 -1191}\special{fp}%
\special{pa 1772 -1230}\special{pa 1772 -1269}\special{fp}\special{pa 1772 -1308}\special{pa 1772 -1347}\special{fp}%
\special{pa 1772 -1386}\special{pa 1772 -1425}\special{fp}\special{pa 1772 -1464}\special{pa 1772 -1503}\special{fp}%
\special{pa 1772 -1542}\special{pa 1772 -1581}\special{fp}\special{pa 1772 -1620}\special{pa 1772 -1659}\special{fp}%
\special{pa 1772 -1698}\special{pa 1772 -1737}\special{fp}\special{pa 1772 -1777}\special{pa 1772 -1816}\special{fp}%
\special{pa 1772 -1855}\special{pa 1772 -1894}\special{fp}\special{pa 1772 -1933}\special{pa 1772 -1972}\special{fp}%
\special{pa 1772 -2011}\special{pa 1772 -2050}\special{fp}\special{pa 1772 -2089}\special{pa 1772 -2128}\special{fp}%
\special{pa 1772 -2167}\special{pa 1772 -2206}\special{fp}\special{pa 1772 -2245}\special{pa 1772 -2284}\special{fp}%
\special{pa 1772 -2323}\special{pa 1772 -2362}\special{fp}%
%
}%
{%
\color[rgb]{0,0,0}%
\special{pa -2362 -1772}\special{pa -2323 -1772}\special{fp}\special{pa -2284 -1772}\special{pa -2245 -1772}\special{fp}%
\special{pa -2206 -1772}\special{pa -2167 -1772}\special{fp}\special{pa -2128 -1772}\special{pa -2089 -1772}\special{fp}%
\special{pa -2050 -1772}\special{pa -2011 -1772}\special{fp}\special{pa -1972 -1772}\special{pa -1933 -1772}\special{fp}%
\special{pa -1894 -1772}\special{pa -1855 -1772}\special{fp}\special{pa -1816 -1772}\special{pa -1777 -1772}\special{fp}%
\special{pa -1737 -1772}\special{pa -1698 -1772}\special{fp}\special{pa -1659 -1772}\special{pa -1620 -1772}\special{fp}%
\special{pa -1581 -1772}\special{pa -1542 -1772}\special{fp}\special{pa -1503 -1772}\special{pa -1464 -1772}\special{fp}%
\special{pa -1425 -1772}\special{pa -1386 -1772}\special{fp}\special{pa -1347 -1772}\special{pa -1308 -1772}\special{fp}%
\special{pa -1269 -1772}\special{pa -1230 -1772}\special{fp}\special{pa -1191 -1772}\special{pa -1152 -1772}\special{fp}%
\special{pa -1113 -1772}\special{pa -1074 -1772}\special{fp}\special{pa -1035 -1772}\special{pa -996 -1772}\special{fp}%
\special{pa -957 -1772}\special{pa -918 -1772}\special{fp}\special{pa -879 -1772}\special{pa -839 -1772}\special{fp}%
\special{pa -800 -1772}\special{pa -761 -1772}\special{fp}\special{pa -722 -1772}\special{pa -683 -1772}\special{fp}%
\special{pa -644 -1772}\special{pa -605 -1772}\special{fp}\special{pa -566 -1772}\special{pa -527 -1772}\special{fp}%
\special{pa -488 -1772}\special{pa -449 -1772}\special{fp}\special{pa -410 -1772}\special{pa -371 -1772}\special{fp}%
\special{pa -332 -1772}\special{pa -293 -1772}\special{fp}\special{pa -254 -1772}\special{pa -215 -1772}\special{fp}%
\special{pa -176 -1772}\special{pa -137 -1772}\special{fp}\special{pa -98 -1772}\special{pa -59 -1772}\special{fp}%
\special{pa -20 -1772}\special{pa 20 -1772}\special{fp}\special{pa 59 -1772}\special{pa 98 -1772}\special{fp}%
\special{pa 137 -1772}\special{pa 176 -1772}\special{fp}\special{pa 215 -1772}\special{pa 254 -1772}\special{fp}%
\special{pa 293 -1772}\special{pa 332 -1772}\special{fp}\special{pa 371 -1772}\special{pa 410 -1772}\special{fp}%
\special{pa 449 -1772}\special{pa 488 -1772}\special{fp}\special{pa 527 -1772}\special{pa 566 -1772}\special{fp}%
\special{pa 605 -1772}\special{pa 644 -1772}\special{fp}\special{pa 683 -1772}\special{pa 722 -1772}\special{fp}%
\special{pa 761 -1772}\special{pa 800 -1772}\special{fp}\special{pa 839 -1772}\special{pa 879 -1772}\special{fp}%
\special{pa 918 -1772}\special{pa 957 -1772}\special{fp}\special{pa 996 -1772}\special{pa 1035 -1772}\special{fp}%
\special{pa 1074 -1772}\special{pa 1113 -1772}\special{fp}\special{pa 1152 -1772}\special{pa 1191 -1772}\special{fp}%
\special{pa 1230 -1772}\special{pa 1269 -1772}\special{fp}\special{pa 1308 -1772}\special{pa 1347 -1772}\special{fp}%
\special{pa 1386 -1772}\special{pa 1425 -1772}\special{fp}\special{pa 1464 -1772}\special{pa 1503 -1772}\special{fp}%
\special{pa 1542 -1772}\special{pa 1581 -1772}\special{fp}\special{pa 1620 -1772}\special{pa 1659 -1772}\special{fp}%
\special{pa 1698 -1772}\special{pa 1737 -1772}\special{fp}\special{pa 1777 -1772}\special{pa 1816 -1772}\special{fp}%
\special{pa 1855 -1772}\special{pa 1894 -1772}\special{fp}\special{pa 1933 -1772}\special{pa 1972 -1772}\special{fp}%
\special{pa 2011 -1772}\special{pa 2050 -1772}\special{fp}\special{pa 2089 -1772}\special{pa 2128 -1772}\special{fp}%
\special{pa 2167 -1772}\special{pa 2206 -1772}\special{fp}\special{pa 2245 -1772}\special{pa 2284 -1772}\special{fp}%
\special{pa 2323 -1772}\special{pa 2362 -1772}\special{fp}%
%
}%
{%
\color[rgb]{0,0,0}%
\special{pa 2165 2362}\special{pa 2165 2323}\special{fp}\special{pa 2165 2284}\special{pa 2165 2245}\special{fp}%
\special{pa 2165 2206}\special{pa 2165 2167}\special{fp}\special{pa 2165 2128}\special{pa 2165 2089}\special{fp}%
\special{pa 2165 2050}\special{pa 2165 2011}\special{fp}\special{pa 2165 1972}\special{pa 2165 1933}\special{fp}%
\special{pa 2165 1894}\special{pa 2165 1855}\special{fp}\special{pa 2165 1816}\special{pa 2165 1777}\special{fp}%
\special{pa 2165 1737}\special{pa 2165 1698}\special{fp}\special{pa 2165 1659}\special{pa 2165 1620}\special{fp}%
\special{pa 2165 1581}\special{pa 2165 1542}\special{fp}\special{pa 2165 1503}\special{pa 2165 1464}\special{fp}%
\special{pa 2165 1425}\special{pa 2165 1386}\special{fp}\special{pa 2165 1347}\special{pa 2165 1308}\special{fp}%
\special{pa 2165 1269}\special{pa 2165 1230}\special{fp}\special{pa 2165 1191}\special{pa 2165 1152}\special{fp}%
\special{pa 2165 1113}\special{pa 2165 1074}\special{fp}\special{pa 2165 1035}\special{pa 2165 996}\special{fp}%
\special{pa 2165 957}\special{pa 2165 918}\special{fp}\special{pa 2165 879}\special{pa 2165 839}\special{fp}%
\special{pa 2165 800}\special{pa 2165 761}\special{fp}\special{pa 2165 722}\special{pa 2165 683}\special{fp}%
\special{pa 2165 644}\special{pa 2165 605}\special{fp}\special{pa 2165 566}\special{pa 2165 527}\special{fp}%
\special{pa 2165 488}\special{pa 2165 449}\special{fp}\special{pa 2165 410}\special{pa 2165 371}\special{fp}%
\special{pa 2165 332}\special{pa 2165 293}\special{fp}\special{pa 2165 254}\special{pa 2165 215}\special{fp}%
\special{pa 2165 176}\special{pa 2165 137}\special{fp}\special{pa 2165 98}\special{pa 2165 59}\special{fp}%
\special{pa 2165 20}\special{pa 2165 -20}\special{fp}\special{pa 2165 -59}\special{pa 2165 -98}\special{fp}%
\special{pa 2165 -137}\special{pa 2165 -176}\special{fp}\special{pa 2165 -215}\special{pa 2165 -254}\special{fp}%
\special{pa 2165 -293}\special{pa 2165 -332}\special{fp}\special{pa 2165 -371}\special{pa 2165 -410}\special{fp}%
\special{pa 2165 -449}\special{pa 2165 -488}\special{fp}\special{pa 2165 -527}\special{pa 2165 -566}\special{fp}%
\special{pa 2165 -605}\special{pa 2165 -644}\special{fp}\special{pa 2165 -683}\special{pa 2165 -722}\special{fp}%
\special{pa 2165 -761}\special{pa 2165 -800}\special{fp}\special{pa 2165 -839}\special{pa 2165 -879}\special{fp}%
\special{pa 2165 -918}\special{pa 2165 -957}\special{fp}\special{pa 2165 -996}\special{pa 2165 -1035}\special{fp}%
\special{pa 2165 -1074}\special{pa 2165 -1113}\special{fp}\special{pa 2165 -1152}\special{pa 2165 -1191}\special{fp}%
\special{pa 2165 -1230}\special{pa 2165 -1269}\special{fp}\special{pa 2165 -1308}\special{pa 2165 -1347}\special{fp}%
\special{pa 2165 -1386}\special{pa 2165 -1425}\special{fp}\special{pa 2165 -1464}\special{pa 2165 -1503}\special{fp}%
\special{pa 2165 -1542}\special{pa 2165 -1581}\special{fp}\special{pa 2165 -1620}\special{pa 2165 -1659}\special{fp}%
\special{pa 2165 -1698}\special{pa 2165 -1737}\special{fp}\special{pa 2165 -1777}\special{pa 2165 -1816}\special{fp}%
\special{pa 2165 -1855}\special{pa 2165 -1894}\special{fp}\special{pa 2165 -1933}\special{pa 2165 -1972}\special{fp}%
\special{pa 2165 -2011}\special{pa 2165 -2050}\special{fp}\special{pa 2165 -2089}\special{pa 2165 -2128}\special{fp}%
\special{pa 2165 -2167}\special{pa 2165 -2206}\special{fp}\special{pa 2165 -2245}\special{pa 2165 -2284}\special{fp}%
\special{pa 2165 -2323}\special{pa 2165 -2362}\special{fp}%
%
}%
{%
\color[rgb]{0,0,0}%
\special{pa -2362 -2165}\special{pa -2323 -2165}\special{fp}\special{pa -2284 -2165}\special{pa -2245 -2165}\special{fp}%
\special{pa -2206 -2165}\special{pa -2167 -2165}\special{fp}\special{pa -2128 -2165}\special{pa -2089 -2165}\special{fp}%
\special{pa -2050 -2165}\special{pa -2011 -2165}\special{fp}\special{pa -1972 -2165}\special{pa -1933 -2165}\special{fp}%
\special{pa -1894 -2165}\special{pa -1855 -2165}\special{fp}\special{pa -1816 -2165}\special{pa -1777 -2165}\special{fp}%
\special{pa -1737 -2165}\special{pa -1698 -2165}\special{fp}\special{pa -1659 -2165}\special{pa -1620 -2165}\special{fp}%
\special{pa -1581 -2165}\special{pa -1542 -2165}\special{fp}\special{pa -1503 -2165}\special{pa -1464 -2165}\special{fp}%
\special{pa -1425 -2165}\special{pa -1386 -2165}\special{fp}\special{pa -1347 -2165}\special{pa -1308 -2165}\special{fp}%
\special{pa -1269 -2165}\special{pa -1230 -2165}\special{fp}\special{pa -1191 -2165}\special{pa -1152 -2165}\special{fp}%
\special{pa -1113 -2165}\special{pa -1074 -2165}\special{fp}\special{pa -1035 -2165}\special{pa -996 -2165}\special{fp}%
\special{pa -957 -2165}\special{pa -918 -2165}\special{fp}\special{pa -879 -2165}\special{pa -839 -2165}\special{fp}%
\special{pa -800 -2165}\special{pa -761 -2165}\special{fp}\special{pa -722 -2165}\special{pa -683 -2165}\special{fp}%
\special{pa -644 -2165}\special{pa -605 -2165}\special{fp}\special{pa -566 -2165}\special{pa -527 -2165}\special{fp}%
\special{pa -488 -2165}\special{pa -449 -2165}\special{fp}\special{pa -410 -2165}\special{pa -371 -2165}\special{fp}%
\special{pa -332 -2165}\special{pa -293 -2165}\special{fp}\special{pa -254 -2165}\special{pa -215 -2165}\special{fp}%
\special{pa -176 -2165}\special{pa -137 -2165}\special{fp}\special{pa -98 -2165}\special{pa -59 -2165}\special{fp}%
\special{pa -20 -2165}\special{pa 20 -2165}\special{fp}\special{pa 59 -2165}\special{pa 98 -2165}\special{fp}%
\special{pa 137 -2165}\special{pa 176 -2165}\special{fp}\special{pa 215 -2165}\special{pa 254 -2165}\special{fp}%
\special{pa 293 -2165}\special{pa 332 -2165}\special{fp}\special{pa 371 -2165}\special{pa 410 -2165}\special{fp}%
\special{pa 449 -2165}\special{pa 488 -2165}\special{fp}\special{pa 527 -2165}\special{pa 566 -2165}\special{fp}%
\special{pa 605 -2165}\special{pa 644 -2165}\special{fp}\special{pa 683 -2165}\special{pa 722 -2165}\special{fp}%
\special{pa 761 -2165}\special{pa 800 -2165}\special{fp}\special{pa 839 -2165}\special{pa 879 -2165}\special{fp}%
\special{pa 918 -2165}\special{pa 957 -2165}\special{fp}\special{pa 996 -2165}\special{pa 1035 -2165}\special{fp}%
\special{pa 1074 -2165}\special{pa 1113 -2165}\special{fp}\special{pa 1152 -2165}\special{pa 1191 -2165}\special{fp}%
\special{pa 1230 -2165}\special{pa 1269 -2165}\special{fp}\special{pa 1308 -2165}\special{pa 1347 -2165}\special{fp}%
\special{pa 1386 -2165}\special{pa 1425 -2165}\special{fp}\special{pa 1464 -2165}\special{pa 1503 -2165}\special{fp}%
\special{pa 1542 -2165}\special{pa 1581 -2165}\special{fp}\special{pa 1620 -2165}\special{pa 1659 -2165}\special{fp}%
\special{pa 1698 -2165}\special{pa 1737 -2165}\special{fp}\special{pa 1777 -2165}\special{pa 1816 -2165}\special{fp}%
\special{pa 1855 -2165}\special{pa 1894 -2165}\special{fp}\special{pa 1933 -2165}\special{pa 1972 -2165}\special{fp}%
\special{pa 2011 -2165}\special{pa 2050 -2165}\special{fp}\special{pa 2089 -2165}\special{pa 2128 -2165}\special{fp}%
\special{pa 2167 -2165}\special{pa 2206 -2165}\special{fp}\special{pa 2245 -2165}\special{pa 2284 -2165}\special{fp}%
\special{pa 2323 -2165}\special{pa 2362 -2165}\special{fp}%
%
}%
\special{pn 8}%
{%
\color[rgb]{0,0,0}%
\special{pa  2362   -20}\special{pa  2362    20}%
\special{fp}%
}%
{%
\color[rgb]{0,0,0}%
\settowidth{\Width}{$-6$}\setlength{\Width}{-0.5\Width}%
\settoheight{\Height}{$-6$}\settodepth{\Depth}{$-6$}\setlength{\Height}{-\Height}%
\put(-6.0000000,-0.1000000){\hspace*{\Width}\raisebox{\Height}{$-6$}}%
%
}%
{%
\color[rgb]{0,0,0}%
\special{pa    20 -2362}\special{pa   -20 -2362}%
\special{fp}%
}%
{%
\color[rgb]{0,0,0}%
\settowidth{\Width}{$-6$}\setlength{\Width}{-1\Width}%
\settoheight{\Height}{$-6$}\settodepth{\Depth}{$-6$}\setlength{\Height}{-0.5\Height}\setlength{\Depth}{0.5\Depth}\addtolength{\Height}{\Depth}%
\put(-0.1000000,-6.0000000){\hspace*{\Width}\raisebox{\Height}{$-6$}}%
%
}%
{%
\color[rgb]{0,0,0}%
\special{pa  2362   -20}\special{pa  2362    20}%
\special{fp}%
}%
{%
\color[rgb]{0,0,0}%
\settowidth{\Width}{$-5$}\setlength{\Width}{-0.5\Width}%
\settoheight{\Height}{$-5$}\settodepth{\Depth}{$-5$}\setlength{\Height}{-\Height}%
\put(-5.0000000,-0.1000000){\hspace*{\Width}\raisebox{\Height}{$-5$}}%
%
}%
{%
\color[rgb]{0,0,0}%
\special{pa    20 -2362}\special{pa   -20 -2362}%
\special{fp}%
}%
{%
\color[rgb]{0,0,0}%
\settowidth{\Width}{$-5$}\setlength{\Width}{-1\Width}%
\settoheight{\Height}{$-5$}\settodepth{\Depth}{$-5$}\setlength{\Height}{-0.5\Height}\setlength{\Depth}{0.5\Depth}\addtolength{\Height}{\Depth}%
\put(-0.1000000,-5.0000000){\hspace*{\Width}\raisebox{\Height}{$-5$}}%
%
}%
{%
\color[rgb]{0,0,0}%
\special{pa  2362   -20}\special{pa  2362    20}%
\special{fp}%
}%
{%
\color[rgb]{0,0,0}%
\settowidth{\Width}{$-4$}\setlength{\Width}{-0.5\Width}%
\settoheight{\Height}{$-4$}\settodepth{\Depth}{$-4$}\setlength{\Height}{-\Height}%
\put(-4.0000000,-0.1000000){\hspace*{\Width}\raisebox{\Height}{$-4$}}%
%
}%
{%
\color[rgb]{0,0,0}%
\special{pa    20 -2362}\special{pa   -20 -2362}%
\special{fp}%
}%
{%
\color[rgb]{0,0,0}%
\settowidth{\Width}{$-4$}\setlength{\Width}{-1\Width}%
\settoheight{\Height}{$-4$}\settodepth{\Depth}{$-4$}\setlength{\Height}{-0.5\Height}\setlength{\Depth}{0.5\Depth}\addtolength{\Height}{\Depth}%
\put(-0.1000000,-4.0000000){\hspace*{\Width}\raisebox{\Height}{$-4$}}%
%
}%
{%
\color[rgb]{0,0,0}%
\special{pa  2362   -20}\special{pa  2362    20}%
\special{fp}%
}%
{%
\color[rgb]{0,0,0}%
\settowidth{\Width}{$-3$}\setlength{\Width}{-0.5\Width}%
\settoheight{\Height}{$-3$}\settodepth{\Depth}{$-3$}\setlength{\Height}{-\Height}%
\put(-3.0000000,-0.1000000){\hspace*{\Width}\raisebox{\Height}{$-3$}}%
%
}%
{%
\color[rgb]{0,0,0}%
\special{pa    20 -2362}\special{pa   -20 -2362}%
\special{fp}%
}%
{%
\color[rgb]{0,0,0}%
\settowidth{\Width}{$-3$}\setlength{\Width}{-1\Width}%
\settoheight{\Height}{$-3$}\settodepth{\Depth}{$-3$}\setlength{\Height}{-0.5\Height}\setlength{\Depth}{0.5\Depth}\addtolength{\Height}{\Depth}%
\put(-0.1000000,-3.0000000){\hspace*{\Width}\raisebox{\Height}{$-3$}}%
%
}%
{%
\color[rgb]{0,0,0}%
\special{pa  2362   -20}\special{pa  2362    20}%
\special{fp}%
}%
{%
\color[rgb]{0,0,0}%
\settowidth{\Width}{$-2$}\setlength{\Width}{-0.5\Width}%
\settoheight{\Height}{$-2$}\settodepth{\Depth}{$-2$}\setlength{\Height}{-\Height}%
\put(-2.0000000,-0.1000000){\hspace*{\Width}\raisebox{\Height}{$-2$}}%
%
}%
{%
\color[rgb]{0,0,0}%
\special{pa    20 -2362}\special{pa   -20 -2362}%
\special{fp}%
}%
{%
\color[rgb]{0,0,0}%
\settowidth{\Width}{$-2$}\setlength{\Width}{-1\Width}%
\settoheight{\Height}{$-2$}\settodepth{\Depth}{$-2$}\setlength{\Height}{-0.5\Height}\setlength{\Depth}{0.5\Depth}\addtolength{\Height}{\Depth}%
\put(-0.1000000,-2.0000000){\hspace*{\Width}\raisebox{\Height}{$-2$}}%
%
}%
{%
\color[rgb]{0,0,0}%
\special{pa  2362   -20}\special{pa  2362    20}%
\special{fp}%
}%
{%
\color[rgb]{0,0,0}%
\settowidth{\Width}{$-1$}\setlength{\Width}{-0.5\Width}%
\settoheight{\Height}{$-1$}\settodepth{\Depth}{$-1$}\setlength{\Height}{-\Height}%
\put(-1.0000000,-0.1000000){\hspace*{\Width}\raisebox{\Height}{$-1$}}%
%
}%
{%
\color[rgb]{0,0,0}%
\special{pa    20 -2362}\special{pa   -20 -2362}%
\special{fp}%
}%
{%
\color[rgb]{0,0,0}%
\settowidth{\Width}{$-1$}\setlength{\Width}{-1\Width}%
\settoheight{\Height}{$-1$}\settodepth{\Depth}{$-1$}\setlength{\Height}{-0.5\Height}\setlength{\Depth}{0.5\Depth}\addtolength{\Height}{\Depth}%
\put(-0.1000000,-1.0000000){\hspace*{\Width}\raisebox{\Height}{$-1$}}%
%
}%
{%
\color[rgb]{0,0,0}%
\special{pa  2362   -20}\special{pa  2362    20}%
\special{fp}%
}%
{%
\color[rgb]{0,0,0}%
\settowidth{\Width}{$1$}\setlength{\Width}{-0.5\Width}%
\settoheight{\Height}{$1$}\settodepth{\Depth}{$1$}\setlength{\Height}{-\Height}%
\put(1.0000000,-0.1000000){\hspace*{\Width}\raisebox{\Height}{$1$}}%
%
}%
{%
\color[rgb]{0,0,0}%
\special{pa    20 -2362}\special{pa   -20 -2362}%
\special{fp}%
}%
{%
\color[rgb]{0,0,0}%
\settowidth{\Width}{$1$}\setlength{\Width}{-1\Width}%
\settoheight{\Height}{$1$}\settodepth{\Depth}{$1$}\setlength{\Height}{-0.5\Height}\setlength{\Depth}{0.5\Depth}\addtolength{\Height}{\Depth}%
\put(-0.1000000,1.0000000){\hspace*{\Width}\raisebox{\Height}{$1$}}%
%
}%
{%
\color[rgb]{0,0,0}%
\special{pa  2362   -20}\special{pa  2362    20}%
\special{fp}%
}%
{%
\color[rgb]{0,0,0}%
\settowidth{\Width}{$2$}\setlength{\Width}{-0.5\Width}%
\settoheight{\Height}{$2$}\settodepth{\Depth}{$2$}\setlength{\Height}{-\Height}%
\put(2.0000000,-0.1000000){\hspace*{\Width}\raisebox{\Height}{$2$}}%
%
}%
{%
\color[rgb]{0,0,0}%
\special{pa    20 -2362}\special{pa   -20 -2362}%
\special{fp}%
}%
{%
\color[rgb]{0,0,0}%
\settowidth{\Width}{$2$}\setlength{\Width}{-1\Width}%
\settoheight{\Height}{$2$}\settodepth{\Depth}{$2$}\setlength{\Height}{-0.5\Height}\setlength{\Depth}{0.5\Depth}\addtolength{\Height}{\Depth}%
\put(-0.1000000,2.0000000){\hspace*{\Width}\raisebox{\Height}{$2$}}%
%
}%
{%
\color[rgb]{0,0,0}%
\special{pa  2362   -20}\special{pa  2362    20}%
\special{fp}%
}%
{%
\color[rgb]{0,0,0}%
\settowidth{\Width}{$3$}\setlength{\Width}{-0.5\Width}%
\settoheight{\Height}{$3$}\settodepth{\Depth}{$3$}\setlength{\Height}{-\Height}%
\put(3.0000000,-0.1000000){\hspace*{\Width}\raisebox{\Height}{$3$}}%
%
}%
{%
\color[rgb]{0,0,0}%
\special{pa    20 -2362}\special{pa   -20 -2362}%
\special{fp}%
}%
{%
\color[rgb]{0,0,0}%
\settowidth{\Width}{$3$}\setlength{\Width}{-1\Width}%
\settoheight{\Height}{$3$}\settodepth{\Depth}{$3$}\setlength{\Height}{-0.5\Height}\setlength{\Depth}{0.5\Depth}\addtolength{\Height}{\Depth}%
\put(-0.1000000,3.0000000){\hspace*{\Width}\raisebox{\Height}{$3$}}%
%
}%
{%
\color[rgb]{0,0,0}%
\special{pa  2362   -20}\special{pa  2362    20}%
\special{fp}%
}%
{%
\color[rgb]{0,0,0}%
\settowidth{\Width}{$4$}\setlength{\Width}{-0.5\Width}%
\settoheight{\Height}{$4$}\settodepth{\Depth}{$4$}\setlength{\Height}{-\Height}%
\put(4.0000000,-0.1000000){\hspace*{\Width}\raisebox{\Height}{$4$}}%
%
}%
{%
\color[rgb]{0,0,0}%
\special{pa    20 -2362}\special{pa   -20 -2362}%
\special{fp}%
}%
{%
\color[rgb]{0,0,0}%
\settowidth{\Width}{$4$}\setlength{\Width}{-1\Width}%
\settoheight{\Height}{$4$}\settodepth{\Depth}{$4$}\setlength{\Height}{-0.5\Height}\setlength{\Depth}{0.5\Depth}\addtolength{\Height}{\Depth}%
\put(-0.1000000,4.0000000){\hspace*{\Width}\raisebox{\Height}{$4$}}%
%
}%
{%
\color[rgb]{0,0,0}%
\special{pa  2362   -20}\special{pa  2362    20}%
\special{fp}%
}%
{%
\color[rgb]{0,0,0}%
\settowidth{\Width}{$5$}\setlength{\Width}{-0.5\Width}%
\settoheight{\Height}{$5$}\settodepth{\Depth}{$5$}\setlength{\Height}{-\Height}%
\put(5.0000000,-0.1000000){\hspace*{\Width}\raisebox{\Height}{$5$}}%
%
}%
{%
\color[rgb]{0,0,0}%
\special{pa    20 -2362}\special{pa   -20 -2362}%
\special{fp}%
}%
{%
\color[rgb]{0,0,0}%
\settowidth{\Width}{$5$}\setlength{\Width}{-1\Width}%
\settoheight{\Height}{$5$}\settodepth{\Depth}{$5$}\setlength{\Height}{-0.5\Height}\setlength{\Depth}{0.5\Depth}\addtolength{\Height}{\Depth}%
\put(-0.1000000,5.0000000){\hspace*{\Width}\raisebox{\Height}{$5$}}%
%
}%
{%
\color[rgb]{0,0,0}%
\special{pa  2362   -20}\special{pa  2362    20}%
\special{fp}%
}%
{%
\color[rgb]{0,0,0}%
\settowidth{\Width}{$6$}\setlength{\Width}{-0.5\Width}%
\settoheight{\Height}{$6$}\settodepth{\Depth}{$6$}\setlength{\Height}{-\Height}%
\put(6.0000000,-0.1000000){\hspace*{\Width}\raisebox{\Height}{$6$}}%
%
}%
{%
\color[rgb]{0,0,0}%
\special{pa    20 -2362}\special{pa   -20 -2362}%
\special{fp}%
}%
{%
\color[rgb]{0,0,0}%
\settowidth{\Width}{$6$}\setlength{\Width}{-1\Width}%
\settoheight{\Height}{$6$}\settodepth{\Depth}{$6$}\setlength{\Height}{-0.5\Height}\setlength{\Depth}{0.5\Depth}\addtolength{\Height}{\Depth}%
\put(-0.1000000,6.0000000){\hspace*{\Width}\raisebox{\Height}{$6$}}%
%
{%
\color[cmyk]{0,1,1,0}%
\special{pa -1157 1969}\special{pa -1159 1960}\special{pa -1164 1953}\special{pa -1171 1947}%
\special{pa -1179 1945}\special{pa -1188 1946}\special{pa -1195 1950}\special{pa -1201 1956}%
\special{pa -1204 1964}\special{pa -1204 1973}\special{pa -1201 1981}\special{pa -1195 1987}%
\special{pa -1188 1991}\special{pa -1179 1992}\special{pa -1171 1990}\special{pa -1164 1984}%
\special{pa -1159 1977}\special{pa -1157 1969}\special{pa -1157 1969}\special{sh 1}\special{ip}%
}%
}%
{%
\color[cmyk]{0,1,1,0}%
\special{pa -1157  1969}\special{pa -1159  1960}\special{pa -1164  1953}\special{pa -1171  1947}%
\special{pa -1179  1945}\special{pa -1188  1946}\special{pa -1195  1950}\special{pa -1201  1956}%
\special{pa -1204  1964}\special{pa -1204  1973}\special{pa -1201  1981}\special{pa -1195  1987}%
\special{pa -1188  1991}\special{pa -1179  1992}\special{pa -1171  1990}\special{pa -1164  1984}%
\special{pa -1159  1977}\special{pa -1157  1969}%
\special{fp}%
{%
\color[cmyk]{0,1,1,0}%
\special{pa -764 1181}\special{pa -765 1173}\special{pa -770 1165}\special{pa -777 1160}%
\special{pa -785 1158}\special{pa -794 1158}\special{pa -802 1162}\special{pa -807 1169}%
\special{pa -811 1177}\special{pa -811 1185}\special{pa -807 1194}\special{pa -802 1200}%
\special{pa -794 1204}\special{pa -785 1205}\special{pa -777 1202}\special{pa -770 1197}%
\special{pa -765 1190}\special{pa -764 1181}\special{pa -764 1181}\special{sh 1}\special{ip}%
}%
}%
{%
\color[cmyk]{0,1,1,0}%
\special{pa  -764  1181}\special{pa  -765  1173}\special{pa  -770  1165}\special{pa  -777  1160}%
\special{pa  -785  1158}\special{pa  -794  1158}\special{pa  -802  1162}\special{pa  -807  1169}%
\special{pa  -811  1177}\special{pa  -811  1185}\special{pa  -807  1194}\special{pa  -802  1200}%
\special{pa  -794  1204}\special{pa  -785  1205}\special{pa  -777  1202}\special{pa  -770  1197}%
\special{pa  -765  1190}\special{pa  -764  1181}%
\special{fp}%
{%
\color[cmyk]{0,1,1,0}%
\special{pa -370 394}\special{pa -372 385}\special{pa -376 378}\special{pa -383 373}%
\special{pa -392 370}\special{pa -400 371}\special{pa -408 375}\special{pa -414 381}%
\special{pa -417 389}\special{pa -417 398}\special{pa -414 406}\special{pa -408 413}%
\special{pa -400 416}\special{pa -392 417}\special{pa -383 415}\special{pa -376 410}%
\special{pa -372 402}\special{pa -370 394}\special{pa -370 394}\special{sh 1}\special{ip}%
}%
}%
{%
\color[cmyk]{0,1,1,0}%
\special{pa  -370   394}\special{pa  -372   385}\special{pa  -376   378}\special{pa  -383   373}%
\special{pa  -392   370}\special{pa  -400   371}\special{pa  -408   375}\special{pa  -414   381}%
\special{pa  -417   389}\special{pa  -417   398}\special{pa  -414   406}\special{pa  -408   413}%
\special{pa  -400   416}\special{pa  -392   417}\special{pa  -383   415}\special{pa  -376   410}%
\special{pa  -372   402}\special{pa  -370   394}%
\special{fp}%
{%
\color[cmyk]{0,1,1,0}%
\special{pa 24 -394}\special{pa 22 -402}\special{pa 17 -410}\special{pa 11 -415}\special{pa 2 -417}%
\special{pa -6 -416}\special{pa -14 -413}\special{pa -20 -406}\special{pa -23 -398}%
\special{pa -23 -389}\special{pa -20 -381}\special{pa -14 -375}\special{pa -6 -371}%
\special{pa 2 -370}\special{pa 11 -373}\special{pa 17 -378}\special{pa 22 -385}\special{pa 24 -394}%
\special{pa 24 -394}\special{sh 1}\special{ip}%
}%
}%
{%
\color[cmyk]{0,1,1,0}%
\special{pa    24  -394}\special{pa    22  -402}\special{pa    17  -410}\special{pa    11  -415}%
\special{pa     2  -417}\special{pa    -6  -416}\special{pa   -14  -413}\special{pa   -20  -406}%
\special{pa   -23  -398}\special{pa   -23  -389}\special{pa   -20  -381}\special{pa   -14  -375}%
\special{pa    -6  -371}\special{pa     2  -370}\special{pa    11  -373}\special{pa    17  -378}%
\special{pa    22  -385}\special{pa    24  -394}%
\special{fp}%
{%
\color[cmyk]{0,1,1,0}%
\special{pa 417 -1181}\special{pa 416 -1190}\special{pa 411 -1197}\special{pa 404 -1202}%
\special{pa 396 -1205}\special{pa 387 -1204}\special{pa 379 -1200}\special{pa 374 -1194}%
\special{pa 370 -1185}\special{pa 370 -1177}\special{pa 374 -1169}\special{pa 379 -1162}%
\special{pa 387 -1158}\special{pa 396 -1158}\special{pa 404 -1160}\special{pa 411 -1165}%
\special{pa 416 -1173}\special{pa 417 -1181}\special{pa 417 -1181}\special{sh 1}\special{ip}%
}%
}%
{%
\color[cmyk]{0,1,1,0}%
\special{pa   417 -1181}\special{pa   416 -1190}\special{pa   411 -1197}\special{pa   404 -1202}%
\special{pa   396 -1205}\special{pa   387 -1204}\special{pa   379 -1200}\special{pa   374 -1194}%
\special{pa   370 -1185}\special{pa   370 -1177}\special{pa   374 -1169}\special{pa   379 -1162}%
\special{pa   387 -1158}\special{pa   396 -1158}\special{pa   404 -1160}\special{pa   411 -1165}%
\special{pa   416 -1173}\special{pa   417 -1181}%
\special{fp}%
{%
\color[cmyk]{0,1,1,0}%
\special{pa 811 -1969}\special{pa 809 -1977}\special{pa 805 -1984}\special{pa 798 -1990}%
\special{pa 790 -1992}\special{pa 781 -1991}\special{pa 773 -1987}\special{pa 767 -1981}%
\special{pa 764 -1973}\special{pa 764 -1964}\special{pa 767 -1956}\special{pa 773 -1950}%
\special{pa 781 -1946}\special{pa 790 -1945}\special{pa 798 -1947}\special{pa 805 -1953}%
\special{pa 809 -1960}\special{pa 811 -1969}\special{pa 811 -1969}\special{sh 1}\special{ip}%
}%
}%
{%
\color[cmyk]{0,1,1,0}%
\special{pa   811 -1969}\special{pa   809 -1977}\special{pa   805 -1984}\special{pa   798 -1990}%
\special{pa   790 -1992}\special{pa   781 -1991}\special{pa   773 -1987}\special{pa   767 -1981}%
\special{pa   764 -1973}\special{pa   764 -1964}\special{pa   767 -1956}\special{pa   773 -1950}%
\special{pa   781 -1946}\special{pa   790 -1945}\special{pa   798 -1947}\special{pa   805 -1953}%
\special{pa   809 -1960}\special{pa   811 -1969}%
\special{fp}%
}%
{%
\color[cmyk]{0,1,1,0}%
\special{pa -1378  2362}\special{pa   984 -2362}%
\special{fp}%
}%
\special{pa -2441    -0}\special{pa  2441    -0}%
\special{fp}%
\special{pa     0  2441}\special{pa     0 -2441}%
\special{fp}%
\settowidth{\Width}{$x$}\setlength{\Width}{0\Width}%
\settoheight{\Height}{$x$}\settodepth{\Depth}{$x$}\setlength{\Height}{-0.5\Height}\setlength{\Depth}{0.5\Depth}\addtolength{\Height}{\Depth}%
\put(6.2500000,0.0000000){\hspace*{\Width}\raisebox{\Height}{$x$}}%
%
\settowidth{\Width}{$y$}\setlength{\Width}{-0.5\Width}%
\settoheight{\Height}{$y$}\settodepth{\Depth}{$y$}\setlength{\Height}{\Depth}%
\put(0.0000000,6.2500000){\hspace*{\Width}\raisebox{\Height}{$y$}}%
%
\settowidth{\Width}{O}\setlength{\Width}{-1\Width}%
\settoheight{\Height}{O}\settodepth{\Depth}{O}\setlength{\Height}{-\Height}%
\put(-0.0500000,-0.0500000){\hspace*{\Width}\raisebox{\Height}{O}}%
%
\end{picture}}%}}
\end{layer}


\newslide{1次関数のグラフ}

\vspace*{18mm}

%%repeat=3
\slidepage
資料の「関数のグラフ」で次の1次関数のグラフをかいてみよう．
\begin{itemize}
\item
$y=3x+3$（TextP81）
\item
$y=10-2x$（TextP81）
\item
$y=2x+2$（TextP81）
\item
$y=\dfrac{1}{2}x+1$
\end{itemize}
\fi %from 259
\label{pageend}\mbox{}

\end{document}
